% Figura (simulación): esfera uniformemente polarizada (problema 4.1).
% Datos: scripts/lineas_campo.py -> datos/esfera-polarizada-{E,D}.tex
\begin{figure}[t]
\centering
\pgfplotsset{esferap/.style={marco, width=7.4cm, xmin=-2.6, xmax=2.6, ymin=-2.0, ymax=2.0}}
\begin{subfigure}[b]{0.49\textwidth}\centering
\begin{tikzpicture}
\begin{axis}[esferap]
  \fill[dielec] (axis cs:0,0) circle[radius=1];
  %% Generado por scripts/lineas_campo.py -- no editar a mano
\addplot[equipot] coordinates {(-0.905,-0.429) (-0.926,-0.401) (-0.946,-0.372) (-0.964,-0.342) (-0.981,-0.311) (-0.996,-0.280) (-1.009,-0.247) (-1.021,-0.214) (-1.031,-0.180) (-1.039,-0.146) (-1.045,-0.112) (-1.050,-0.077) (-1.053,-0.042) (-1.054,-0.007) (-1.054,0.028) (-1.051,0.063) (-1.047,0.098) (-1.041,0.133) (-1.034,0.167) (-1.025,0.201) (-1.014,0.234) (-1.001,0.267) (-0.987,0.299) (-0.971,0.330) (-0.953,0.361) (-0.934,0.390) (-0.914,0.419) (-0.900,0.411) (-0.900,0.375) (-0.900,0.340) (-0.900,0.305) (-0.900,0.270) (-0.900,0.235) (-0.900,0.200) (-0.900,0.165) (-0.900,0.130) (-0.900,0.095) (-0.900,0.059) (-0.900,0.024) (-0.900,-0.011) (-0.900,-0.046) (-0.900,-0.081) (-0.900,-0.116) (-0.900,-0.151) (-0.900,-0.186) (-0.900,-0.221) (-0.900,-0.257) (-0.900,-0.292) (-0.900,-0.327) (-0.900,-0.362) (-0.900,-0.397) (-0.905,-0.429)};
\addplot[equipot] coordinates {(-0.726,-0.692) (-0.759,-0.678) (-0.790,-0.663) (-0.821,-0.646) (-0.852,-0.627) (-0.881,-0.607) (-0.908,-0.585) (-0.935,-0.562) (-0.961,-0.538) (-0.985,-0.512) (-1.008,-0.485) (-1.030,-0.457) (-1.050,-0.428) (-1.069,-0.398) (-1.087,-0.367) (-1.103,-0.336) (-1.117,-0.304) (-1.130,-0.271) (-1.142,-0.237) (-1.152,-0.203) (-1.160,-0.169) (-1.167,-0.134) (-1.172,-0.099) (-1.176,-0.064) (-1.178,-0.029) (-1.178,0.007) (-1.177,0.042) (-1.175,0.077) (-1.170,0.112) (-1.165,0.147) (-1.157,0.182) (-1.148,0.216) (-1.138,0.250) (-1.126,0.283) (-1.112,0.316) (-1.097,0.348) (-1.080,0.379) (-1.062,0.409) (-1.043,0.439) (-1.022,0.468) (-1.000,0.495) (-0.976,0.522) (-0.951,0.547) (-0.925,0.571) (-0.898,0.593) (-0.870,0.615) (-0.840,0.634) (-0.810,0.652) (-0.779,0.669) (-0.746,0.683) (-0.720,0.683) (-0.720,0.648) (-0.720,0.613) (-0.720,0.577) (-0.720,0.542) (-0.720,0.506) (-0.720,0.471) (-0.720,0.436) (-0.720,0.400) (-0.720,0.365) (-0.720,0.330) (-0.720,0.294) (-0.720,0.259) (-0.720,0.223) (-0.720,0.188) (-0.720,0.153) (-0.720,0.117) (-0.720,0.082) (-0.720,0.046) (-0.720,0.011) (-0.720,-0.024) (-0.720,-0.060) (-0.720,-0.095) (-0.720,-0.130) (-0.720,-0.166) (-0.720,-0.201) (-0.720,-0.237) (-0.720,-0.272) (-0.720,-0.307) (-0.720,-0.343) (-0.720,-0.378) (-0.720,-0.414) (-0.720,-0.449) (-0.720,-0.484) (-0.720,-0.520) (-0.720,-0.555) (-0.720,-0.590) (-0.720,-0.626) (-0.720,-0.661) (-0.726,-0.692)};
\addplot[equipot] coordinates {(-0.662,-0.841) (-0.697,-0.836) (-0.732,-0.830) (-0.766,-0.822) (-0.799,-0.812) (-0.833,-0.801) (-0.865,-0.788) (-0.897,-0.773) (-0.928,-0.757) (-0.959,-0.740) (-0.988,-0.721) (-1.017,-0.700) (-1.045,-0.679) (-1.072,-0.656) (-1.097,-0.633) (-1.122,-0.608) (-1.146,-0.582) (-1.168,-0.555) (-1.190,-0.527) (-1.210,-0.498) (-1.229,-0.469) (-1.247,-0.439) (-1.263,-0.408) (-1.279,-0.376) (-1.293,-0.344) (-1.306,-0.311) (-1.317,-0.278) (-1.327,-0.244) (-1.336,-0.210) (-1.344,-0.176) (-1.350,-0.142) (-1.355,-0.107) (-1.358,-0.072) (-1.360,-0.037) (-1.361,-0.002) (-1.360,0.033) (-1.358,0.068) (-1.355,0.103) (-1.350,0.138) (-1.344,0.173) (-1.337,0.207) (-1.328,0.241) (-1.318,0.275) (-1.307,0.308) (-1.294,0.341) (-1.280,0.373) (-1.265,0.404) (-1.249,0.435) (-1.231,0.466) (-1.212,0.495) (-1.192,0.524) (-1.171,0.552) (-1.148,0.579) (-1.125,0.605) (-1.100,0.630) (-1.074,0.654) (-1.048,0.677) (-1.020,0.698) (-0.991,0.719) (-0.962,0.738) (-0.932,0.755) (-0.900,0.772) (-0.869,0.786) (-0.836,0.800) (-0.803,0.811) (-0.769,0.821) (-0.735,0.829) (-0.701,0.836) (-0.666,0.841) (-0.631,0.843) (-0.596,0.844) (-0.561,0.843) (-0.540,0.827) (-0.540,0.792) (-0.540,0.756) (-0.540,0.721) (-0.540,0.686) (-0.540,0.651) (-0.540,0.616) (-0.540,0.581) (-0.540,0.546) (-0.540,0.511) (-0.540,0.476) (-0.540,0.441) (-0.540,0.405) (-0.540,0.370) (-0.540,0.335) (-0.540,0.300) (-0.540,0.265) (-0.540,0.230) (-0.540,0.195) (-0.540,0.160) (-0.540,0.125) (-0.540,0.090) (-0.540,0.054) (-0.540,0.019) (-0.540,-0.016) (-0.540,-0.051) (-0.540,-0.086) (-0.540,-0.121) (-0.540,-0.156) (-0.540,-0.191) (-0.540,-0.226) (-0.540,-0.261) (-0.540,-0.297) (-0.540,-0.332) (-0.540,-0.367) (-0.540,-0.402) (-0.540,-0.437) (-0.540,-0.472) (-0.540,-0.507) (-0.540,-0.542) (-0.540,-0.577) (-0.540,-0.613) (-0.540,-0.648) (-0.540,-0.683) (-0.540,-0.718) (-0.540,-0.753) (-0.540,-0.788) (-0.540,-0.823) (-0.557,-0.843) (-0.592,-0.844) (-0.627,-0.844) (-0.662,-0.841)};
\addplot[equipot] coordinates {(-0.732,-1.034) (-0.767,-1.033) (-0.802,-1.031) (-0.837,-1.027) (-0.872,-1.021) (-0.906,-1.015) (-0.940,-1.006) (-0.974,-0.996) (-1.007,-0.985) (-1.040,-0.973) (-1.073,-0.959) (-1.104,-0.944) (-1.136,-0.928) (-1.166,-0.911) (-1.196,-0.892) (-1.225,-0.873) (-1.254,-0.852) (-1.281,-0.830) (-1.308,-0.808) (-1.334,-0.784) (-1.360,-0.759) (-1.384,-0.734) (-1.407,-0.708) (-1.430,-0.681) (-1.451,-0.653) (-1.472,-0.625) (-1.492,-0.595) (-1.510,-0.566) (-1.528,-0.535) (-1.545,-0.504) (-1.560,-0.473) (-1.575,-0.441) (-1.589,-0.408) (-1.601,-0.375) (-1.612,-0.342) (-1.623,-0.309) (-1.632,-0.275) (-1.640,-0.241) (-1.647,-0.206) (-1.653,-0.171) (-1.658,-0.137) (-1.662,-0.102) (-1.665,-0.067) (-1.666,-0.031) (-1.667,0.004) (-1.666,0.039) (-1.664,0.074) (-1.661,0.109) (-1.657,0.144) (-1.652,0.179) (-1.646,0.213) (-1.639,0.248) (-1.630,0.282) (-1.621,0.316) (-1.610,0.349) (-1.599,0.382) (-1.586,0.415) (-1.572,0.448) (-1.557,0.479) (-1.541,0.511) (-1.524,0.542) (-1.507,0.572) (-1.488,0.602) (-1.468,0.631) (-1.447,0.659) (-1.425,0.687) (-1.402,0.713) (-1.379,0.739) (-1.354,0.765) (-1.329,0.789) (-1.303,0.812) (-1.276,0.835) (-1.248,0.856) (-1.219,0.877) (-1.190,0.896) (-1.160,0.914) (-1.129,0.932) (-1.098,0.947) (-1.066,0.962) (-1.033,0.976) (-1.000,0.988) (-0.967,0.999) (-0.933,1.008) (-0.899,1.016) (-0.864,1.023) (-0.830,1.028) (-0.795,1.031) (-0.759,1.033) (-0.724,1.034) (-0.689,1.033) (-0.654,1.030) (-0.619,1.026) (-0.585,1.020) (-0.550,1.012) (-0.516,1.003) (-0.483,0.992) (-0.450,0.979) (-0.418,0.964) (-0.387,0.948) (-0.360,0.928) (-0.360,0.893) (-0.360,0.858) (-0.360,0.823) (-0.360,0.787) (-0.360,0.752) (-0.360,0.717) (-0.360,0.682) (-0.360,0.647) (-0.360,0.612) (-0.360,0.576) (-0.360,0.541) (-0.360,0.506) (-0.360,0.471) (-0.360,0.436) (-0.360,0.401) (-0.360,0.366) (-0.360,0.330) (-0.360,0.295) (-0.360,0.260) (-0.360,0.225) (-0.360,0.190) (-0.360,0.155) (-0.360,0.119) (-0.360,0.084) (-0.360,0.049) (-0.360,0.014) (-0.360,-0.021) (-0.360,-0.056) (-0.360,-0.092) (-0.360,-0.127) (-0.360,-0.162) (-0.360,-0.197) (-0.360,-0.232) (-0.360,-0.267) (-0.360,-0.303) (-0.360,-0.338) (-0.360,-0.373) (-0.360,-0.408) (-0.360,-0.443) (-0.360,-0.478) (-0.360,-0.514) (-0.360,-0.549) (-0.360,-0.584) (-0.360,-0.619) (-0.360,-0.654) (-0.360,-0.689) (-0.360,-0.725) (-0.360,-0.760) (-0.360,-0.795) (-0.360,-0.830) (-0.360,-0.865) (-0.360,-0.900) (-0.363,-0.935) (-0.393,-0.952) (-0.425,-0.968) (-0.457,-0.982) (-0.490,-0.994) (-0.523,-1.005) (-0.557,-1.014) (-0.592,-1.021) (-0.627,-1.027) (-0.661,-1.031) (-0.697,-1.033) (-0.732,-1.034)};
\addplot[equipot] coordinates {(-1.148,-1.456) (-1.183,-1.452) (-1.218,-1.447) (-1.252,-1.441) (-1.287,-1.434) (-1.321,-1.425) (-1.355,-1.416) (-1.388,-1.406) (-1.421,-1.395) (-1.454,-1.383) (-1.487,-1.370) (-1.519,-1.356) (-1.551,-1.341) (-1.582,-1.325) (-1.613,-1.309) (-1.644,-1.291) (-1.674,-1.273) (-1.703,-1.254) (-1.732,-1.234) (-1.761,-1.214) (-1.789,-1.193) (-1.816,-1.171) (-1.843,-1.148) (-1.869,-1.125) (-1.895,-1.101) (-1.920,-1.077) (-1.945,-1.051) (-1.968,-1.026) (-1.992,-0.999) (-2.014,-0.973) (-2.036,-0.945) (-2.057,-0.917) (-2.078,-0.889) (-2.098,-0.860) (-2.117,-0.831) (-2.136,-0.801) (-2.153,-0.771) (-2.170,-0.740) (-2.187,-0.709) (-2.203,-0.677) (-2.217,-0.646) (-2.232,-0.614) (-2.245,-0.581) (-2.258,-0.549) (-2.270,-0.516) (-2.281,-0.482) (-2.291,-0.449) (-2.301,-0.415) (-2.310,-0.381) (-2.318,-0.347) (-2.326,-0.313) (-2.332,-0.278) (-2.338,-0.244) (-2.343,-0.209) (-2.347,-0.174) (-2.351,-0.139) (-2.354,-0.104) (-2.356,-0.069) (-2.357,-0.034) (-2.357,0.001) (-2.357,0.036) (-2.355,0.071) (-2.353,0.106) (-2.351,0.141) (-2.347,0.176) (-2.343,0.211) (-2.338,0.246) (-2.332,0.280) (-2.325,0.315) (-2.318,0.349) (-2.310,0.383) (-2.301,0.417) (-2.291,0.451) (-2.280,0.484) (-2.269,0.517) (-2.257,0.550) (-2.244,0.583) (-2.231,0.615) (-2.217,0.647) (-2.202,0.679) (-2.186,0.711) (-2.170,0.742) (-2.152,0.772) (-2.135,0.802) (-2.116,0.832) (-2.097,0.862) (-2.077,0.890) (-2.056,0.919) (-2.035,0.947) (-2.013,0.974) (-1.990,1.001) (-1.967,1.027) (-1.943,1.053) (-1.919,1.078) (-1.894,1.102) (-1.868,1.126) (-1.842,1.150) (-1.815,1.172) (-1.787,1.194) (-1.759,1.215) (-1.731,1.235) (-1.702,1.255) (-1.672,1.274) (-1.642,1.292) (-1.611,1.309) (-1.580,1.326) (-1.549,1.342) (-1.517,1.356) (-1.485,1.370) (-1.452,1.383) (-1.420,1.395) (-1.386,1.406) (-1.353,1.417) (-1.319,1.426) (-1.285,1.434) (-1.250,1.441) (-1.216,1.447) (-1.181,1.453) (-1.146,1.457) (-1.111,1.460) (-1.076,1.461) (-1.041,1.462) (-1.006,1.462) (-0.971,1.460) (-0.936,1.458) (-0.901,1.454) (-0.866,1.449) (-0.832,1.443) (-0.798,1.436) (-0.763,1.427) (-0.730,1.418) (-0.696,1.407) (-0.663,1.395) (-0.631,1.382) (-0.599,1.367) (-0.567,1.352) (-0.536,1.335) (-0.506,1.318) (-0.477,1.299) (-0.448,1.279) (-0.420,1.258) (-0.392,1.236) (-0.366,1.213) (-0.340,1.189) (-0.316,1.164) (-0.292,1.138) (-0.269,1.111) (-0.247,1.084) (-0.226,1.056) (-0.207,1.027) (-0.188,0.997) (-0.180,0.964) (-0.180,0.929) (-0.180,0.894) (-0.180,0.859) (-0.180,0.824) (-0.180,0.789) (-0.180,0.753) (-0.180,0.718) (-0.180,0.683) (-0.180,0.648) (-0.180,0.613) (-0.180,0.578) (-0.180,0.543) (-0.180,0.508) (-0.180,0.473) (-0.180,0.438) (-0.180,0.403) (-0.180,0.367) (-0.180,0.332) (-0.180,0.297) (-0.180,0.262) (-0.180,0.227) (-0.180,0.192) (-0.180,0.157) (-0.180,0.122) (-0.180,0.087) (-0.180,0.052) (-0.180,0.017) (-0.180,-0.019) (-0.180,-0.054) (-0.180,-0.089) (-0.180,-0.124) (-0.180,-0.159) (-0.180,-0.194) (-0.180,-0.229) (-0.180,-0.264) (-0.180,-0.299) (-0.180,-0.334) (-0.180,-0.369) (-0.180,-0.404) (-0.180,-0.440) (-0.180,-0.475) (-0.180,-0.510) (-0.180,-0.545) (-0.180,-0.580) (-0.180,-0.615) (-0.180,-0.650) (-0.180,-0.685) (-0.180,-0.720) (-0.180,-0.755) (-0.180,-0.790) (-0.180,-0.826) (-0.180,-0.861) (-0.180,-0.896) (-0.180,-0.931) (-0.180,-0.966) (-0.189,-0.999) (-0.208,-1.028) (-0.228,-1.057) (-0.248,-1.085) (-0.270,-1.113) (-0.293,-1.139) (-0.317,-1.165) (-0.342,-1.190) (-0.367,-1.214) (-0.394,-1.237) (-0.421,-1.259) (-0.449,-1.280) (-0.478,-1.300) (-0.508,-1.319) (-0.538,-1.336) (-0.569,-1.353) (-0.601,-1.368) (-0.633,-1.383) (-0.665,-1.396) (-0.698,-1.408) (-0.732,-1.418) (-0.765,-1.428) (-0.799,-1.436) (-0.834,-1.443) (-0.868,-1.449) (-0.903,-1.454) (-0.938,-1.458) (-0.973,-1.461) (-1.008,-1.462) (-1.043,-1.462) (-1.078,-1.461) (-1.113,-1.459) (-1.148,-1.456)};
\addplot[equipot] coordinates {(0.923,-1.457) (0.958,-1.460) (0.993,-1.462) (1.028,-1.462) (1.063,-1.462) (1.098,-1.460) (1.133,-1.458) (1.168,-1.454) (1.202,-1.450) (1.237,-1.444) (1.271,-1.437) (1.306,-1.429) (1.340,-1.420) (1.373,-1.411) (1.407,-1.400) (1.440,-1.388) (1.473,-1.375) (1.505,-1.362) (1.537,-1.347) (1.568,-1.332) (1.600,-1.316) (1.630,-1.299) (1.661,-1.281) (1.690,-1.263) (1.720,-1.243) (1.748,-1.223) (1.776,-1.202) (1.804,-1.181) (1.831,-1.158) (1.858,-1.135) (1.884,-1.112) (1.909,-1.087) (1.934,-1.063) (1.958,-1.037) (1.981,-1.011) (2.004,-0.984) (2.027,-0.957) (2.048,-0.930) (2.069,-0.901) (2.089,-0.873) (2.109,-0.844) (2.127,-0.814) (2.146,-0.784) (2.163,-0.753) (2.180,-0.723) (2.196,-0.691) (2.211,-0.660) (2.225,-0.628) (2.239,-0.596) (2.252,-0.563) (2.265,-0.530) (2.276,-0.497) (2.287,-0.464) (2.297,-0.430) (2.306,-0.396) (2.315,-0.362) (2.322,-0.328) (2.329,-0.293) (2.336,-0.259) (2.341,-0.224) (2.346,-0.189) (2.349,-0.155) (2.352,-0.120) (2.355,-0.085) (2.356,-0.050) (2.357,-0.014) (2.357,0.021) (2.356,0.056) (2.354,0.091) (2.352,0.126) (2.349,0.161) (2.345,0.196) (2.340,0.230) (2.335,0.265) (2.328,0.299) (2.321,0.334) (2.313,0.368) (2.305,0.402) (2.295,0.436) (2.285,0.469) (2.274,0.503) (2.262,0.536) (2.250,0.569) (2.237,0.601) (2.223,0.633) (2.208,0.665) (2.193,0.697) (2.177,0.728) (2.160,0.759) (2.143,0.789) (2.124,0.819) (2.105,0.849) (2.086,0.878) (2.065,0.906) (2.044,0.934) (2.023,0.962) (2.000,0.989) (1.977,1.016) (1.954,1.042) (1.930,1.067) (1.905,1.092) (1.879,1.116) (1.853,1.139) (1.827,1.162) (1.799,1.184) (1.772,1.206) (1.743,1.227) (1.714,1.247) (1.685,1.266) (1.655,1.284) (1.625,1.302) (1.594,1.319) (1.563,1.335) (1.531,1.350) (1.499,1.364) (1.467,1.378) (1.434,1.390) (1.401,1.402) (1.367,1.412) (1.334,1.422) (1.300,1.431) (1.266,1.438) (1.231,1.445) (1.196,1.450) (1.162,1.455) (1.127,1.458) (1.092,1.461) (1.057,1.462) (1.022,1.462) (0.986,1.461) (0.951,1.459) (0.917,1.456) (0.882,1.452) (0.847,1.446) (0.813,1.439) (0.778,1.431) (0.745,1.422) (0.711,1.412) (0.678,1.400) (0.645,1.388) (0.613,1.374) (0.581,1.359) (0.550,1.343) (0.519,1.326) (0.490,1.307) (0.460,1.288) (0.432,1.267) (0.404,1.246) (0.377,1.223) (0.351,1.199) (0.326,1.175) (0.302,1.149) (0.279,1.123) (0.257,1.096) (0.236,1.068) (0.215,1.039) (0.196,1.010) (0.180,0.979) (0.180,0.944) (0.180,0.909) (0.180,0.874) (0.180,0.839) (0.180,0.804) (0.180,0.769) (0.180,0.734) (0.180,0.699) (0.180,0.664) (0.180,0.629) (0.180,0.593) (0.180,0.558) (0.180,0.523) (0.180,0.488) (0.180,0.453) (0.180,0.418) (0.180,0.383) (0.180,0.348) (0.180,0.313) (0.180,0.278) (0.180,0.243) (0.180,0.207) (0.180,0.172) (0.180,0.137) (0.180,0.102) (0.180,0.067) (0.180,0.032) (0.180,-0.003) (0.180,-0.038) (0.180,-0.073) (0.180,-0.108) (0.180,-0.143) (0.180,-0.178) (0.180,-0.214) (0.180,-0.249) (0.180,-0.284) (0.180,-0.319) (0.180,-0.354) (0.180,-0.389) (0.180,-0.424) (0.180,-0.459) (0.180,-0.494) (0.180,-0.529) (0.180,-0.564) (0.180,-0.600) (0.180,-0.635) (0.180,-0.670) (0.180,-0.705) (0.180,-0.740) (0.180,-0.775) (0.180,-0.810) (0.180,-0.845) (0.180,-0.880) (0.180,-0.915) (0.180,-0.950) (0.181,-0.985) (0.199,-1.015) (0.219,-1.044) (0.239,-1.073) (0.261,-1.101) (0.283,-1.128) (0.306,-1.154) (0.331,-1.179) (0.356,-1.204) (0.382,-1.227) (0.409,-1.249) (0.437,-1.271) (0.465,-1.291) (0.495,-1.310) (0.525,-1.329) (0.555,-1.346) (0.587,-1.362) (0.618,-1.376) (0.651,-1.390) (0.684,-1.402) (0.717,-1.414) (0.750,-1.424) (0.784,-1.433) (0.819,-1.440) (0.853,-1.447) (0.888,-1.452) (0.923,-1.457)};
\addplot[equipot] coordinates {(0.732,-1.034) (0.767,-1.033) (0.802,-1.031) (0.837,-1.027) (0.872,-1.021) (0.906,-1.015) (0.940,-1.006) (0.974,-0.996) (1.007,-0.985) (1.040,-0.973) (1.073,-0.959) (1.104,-0.944) (1.136,-0.928) (1.166,-0.911) (1.196,-0.892) (1.225,-0.873) (1.254,-0.852) (1.281,-0.830) (1.308,-0.808) (1.334,-0.784) (1.360,-0.759) (1.384,-0.734) (1.407,-0.708) (1.430,-0.681) (1.451,-0.653) (1.472,-0.625) (1.492,-0.595) (1.510,-0.566) (1.528,-0.535) (1.545,-0.504) (1.560,-0.473) (1.575,-0.441) (1.589,-0.408) (1.601,-0.375) (1.612,-0.342) (1.623,-0.309) (1.632,-0.275) (1.640,-0.241) (1.647,-0.206) (1.653,-0.171) (1.658,-0.137) (1.662,-0.102) (1.665,-0.067) (1.666,-0.031) (1.667,0.004) (1.666,0.039) (1.664,0.074) (1.661,0.109) (1.657,0.144) (1.652,0.179) (1.646,0.213) (1.639,0.248) (1.630,0.282) (1.621,0.316) (1.610,0.349) (1.599,0.382) (1.586,0.415) (1.572,0.448) (1.557,0.479) (1.541,0.511) (1.524,0.542) (1.507,0.572) (1.488,0.602) (1.468,0.631) (1.447,0.659) (1.425,0.687) (1.402,0.713) (1.379,0.739) (1.354,0.765) (1.329,0.789) (1.303,0.812) (1.276,0.835) (1.248,0.856) (1.219,0.877) (1.190,0.896) (1.160,0.914) (1.129,0.932) (1.098,0.947) (1.066,0.962) (1.033,0.976) (1.000,0.988) (0.967,0.999) (0.933,1.008) (0.899,1.016) (0.864,1.023) (0.830,1.028) (0.795,1.031) (0.759,1.033) (0.724,1.034) (0.689,1.033) (0.654,1.030) (0.619,1.026) (0.585,1.020) (0.550,1.012) (0.516,1.003) (0.483,0.992) (0.450,0.979) (0.418,0.964) (0.387,0.948) (0.360,0.928) (0.360,0.893) (0.360,0.858) (0.360,0.823) (0.360,0.787) (0.360,0.752) (0.360,0.717) (0.360,0.682) (0.360,0.647) (0.360,0.612) (0.360,0.576) (0.360,0.541) (0.360,0.506) (0.360,0.471) (0.360,0.436) (0.360,0.401) (0.360,0.366) (0.360,0.330) (0.360,0.295) (0.360,0.260) (0.360,0.225) (0.360,0.190) (0.360,0.155) (0.360,0.119) (0.360,0.084) (0.360,0.049) (0.360,0.014) (0.360,-0.021) (0.360,-0.056) (0.360,-0.092) (0.360,-0.127) (0.360,-0.162) (0.360,-0.197) (0.360,-0.232) (0.360,-0.267) (0.360,-0.303) (0.360,-0.338) (0.360,-0.373) (0.360,-0.408) (0.360,-0.443) (0.360,-0.478) (0.360,-0.514) (0.360,-0.549) (0.360,-0.584) (0.360,-0.619) (0.360,-0.654) (0.360,-0.689) (0.360,-0.725) (0.360,-0.760) (0.360,-0.795) (0.360,-0.830) (0.360,-0.865) (0.360,-0.900) (0.363,-0.935) (0.393,-0.952) (0.425,-0.968) (0.457,-0.982) (0.490,-0.994) (0.523,-1.005) (0.557,-1.014) (0.592,-1.021) (0.627,-1.027) (0.661,-1.031) (0.697,-1.033) (0.732,-1.034)};
\addplot[equipot] coordinates {(0.541,-0.841) (0.576,-0.844) (0.611,-0.844) (0.646,-0.842) (0.681,-0.839) (0.716,-0.833) (0.750,-0.826) (0.784,-0.817) (0.817,-0.806) (0.850,-0.794) (0.882,-0.780) (0.914,-0.765) (0.945,-0.748) (0.975,-0.730) (1.004,-0.710) (1.032,-0.689) (1.059,-0.667) (1.086,-0.644) (1.111,-0.619) (1.135,-0.594) (1.158,-0.567) (1.180,-0.540) (1.201,-0.512) (1.220,-0.483) (1.239,-0.453) (1.256,-0.422) (1.272,-0.391) (1.286,-0.359) (1.300,-0.327) (1.312,-0.294) (1.323,-0.260) (1.332,-0.226) (1.340,-0.192) (1.347,-0.158) (1.352,-0.123) (1.357,-0.088) (1.359,-0.053) (1.361,-0.018) (1.361,0.017) (1.359,0.052) (1.357,0.087) (1.353,0.122) (1.347,0.157) (1.340,0.191) (1.332,0.225) (1.323,0.259) (1.312,0.292) (1.300,0.325) (1.287,0.358) (1.272,0.390) (1.256,0.421) (1.239,0.452) (1.221,0.482) (1.201,0.511) (1.181,0.539) (1.159,0.567) (1.136,0.593) (1.112,0.619) (1.086,0.643) (1.060,0.666) (1.033,0.688) (1.005,0.709) (0.976,0.729) (0.946,0.747) (0.915,0.764) (0.884,0.780) (0.851,0.794) (0.818,0.806) (0.785,0.817) (0.751,0.826) (0.717,0.833) (0.682,0.839) (0.647,0.842) (0.612,0.844) (0.577,0.844) (0.542,0.842) (0.540,0.808) (0.540,0.773) (0.540,0.738) (0.540,0.703) (0.540,0.668) (0.540,0.632) (0.540,0.597) (0.540,0.562) (0.540,0.527) (0.540,0.492) (0.540,0.457) (0.540,0.422) (0.540,0.387) (0.540,0.352) (0.540,0.317) (0.540,0.281) (0.540,0.246) (0.540,0.211) (0.540,0.176) (0.540,0.141) (0.540,0.106) (0.540,0.071) (0.540,0.036) (0.540,0.001) (0.540,-0.034) (0.540,-0.070) (0.540,-0.105) (0.540,-0.140) (0.540,-0.175) (0.540,-0.210) (0.540,-0.245) (0.540,-0.280) (0.540,-0.315) (0.540,-0.350) (0.540,-0.385) (0.540,-0.421) (0.540,-0.456) (0.540,-0.491) (0.540,-0.526) (0.540,-0.561) (0.540,-0.596) (0.540,-0.631) (0.540,-0.666) (0.540,-0.701) (0.540,-0.737) (0.540,-0.772) (0.540,-0.807) (0.541,-0.841)};
\addplot[equipot] coordinates {(0.720,-0.691) (0.753,-0.681) (0.785,-0.666) (0.816,-0.649) (0.847,-0.630) (0.876,-0.610) (0.904,-0.589) (0.931,-0.566) (0.957,-0.542) (0.981,-0.516) (1.004,-0.490) (1.026,-0.462) (1.047,-0.433) (1.066,-0.403) (1.084,-0.373) (1.100,-0.341) (1.115,-0.309) (1.128,-0.276) (1.140,-0.243) (1.150,-0.209) (1.159,-0.175) (1.166,-0.140) (1.171,-0.105) (1.175,-0.070) (1.178,-0.035) (1.179,0.001) (1.178,0.036) (1.175,0.071) (1.171,0.107) (1.166,0.141) (1.158,0.176) (1.150,0.210) (1.139,0.244) (1.128,0.278) (1.114,0.310) (1.099,0.342) (1.083,0.374) (1.065,0.404) (1.046,0.434) (1.025,0.463) (1.003,0.491) (0.980,0.517) (0.956,0.543) (0.930,0.567) (0.903,0.590) (0.875,0.611) (0.845,0.631) (0.815,0.649) (0.784,0.666) (0.752,0.681) (0.720,0.689) (0.720,0.654) (0.720,0.619) (0.720,0.583) (0.720,0.548) (0.720,0.512) (0.720,0.477) (0.720,0.442) (0.720,0.406) (0.720,0.371) (0.720,0.335) (0.720,0.300) (0.720,0.265) (0.720,0.229) (0.720,0.194) (0.720,0.159) (0.720,0.123) (0.720,0.088) (0.720,0.052) (0.720,0.017) (0.720,-0.018) (0.720,-0.054) (0.720,-0.089) (0.720,-0.125) (0.720,-0.160) (0.720,-0.195) (0.720,-0.231) (0.720,-0.266) (0.720,-0.301) (0.720,-0.337) (0.720,-0.372) (0.720,-0.408) (0.720,-0.443) (0.720,-0.478) (0.720,-0.514) (0.720,-0.549) (0.720,-0.585) (0.720,-0.620) (0.720,-0.655) (0.720,-0.691)};
\addplot[equipot] coordinates {(0.905,-0.429) (0.926,-0.401) (0.946,-0.372) (0.964,-0.342) (0.981,-0.311) (0.996,-0.280) (1.009,-0.247) (1.021,-0.214) (1.031,-0.180) (1.039,-0.146) (1.045,-0.112) (1.050,-0.077) (1.053,-0.042) (1.054,-0.007) (1.054,0.028) (1.051,0.063) (1.047,0.098) (1.041,0.133) (1.034,0.167) (1.025,0.201) (1.014,0.234) (1.001,0.267) (0.987,0.299) (0.971,0.330) (0.953,0.361) (0.934,0.390) (0.914,0.419) (0.900,0.411) (0.900,0.375) (0.900,0.340) (0.900,0.305) (0.900,0.270) (0.900,0.235) (0.900,0.200) (0.900,0.165) (0.900,0.130) (0.900,0.095) (0.900,0.059) (0.900,0.024) (0.900,-0.011) (0.900,-0.046) (0.900,-0.081) (0.900,-0.116) (0.900,-0.151) (0.900,-0.186) (0.900,-0.221) (0.900,-0.257) (0.900,-0.292) (0.900,-0.327) (0.900,-0.362) (0.900,-0.397) (0.905,-0.429)};
\addplot[lineaE, flechaen=0.711] coordinates {(0.975,0.221) (1.009,0.232) (1.043,0.244) (1.076,0.256) (1.110,0.268) (1.143,0.281) (1.176,0.294) (1.210,0.307) (1.243,0.320) (1.276,0.333) (1.309,0.347) (1.342,0.360) (1.375,0.374) (1.407,0.388) (1.440,0.402) (1.473,0.417) (1.505,0.431) (1.538,0.446) (1.570,0.461) (1.603,0.476) (1.635,0.491) (1.667,0.506) (1.699,0.521) (1.732,0.537) (1.764,0.553) (1.795,0.569) (1.827,0.585) (1.859,0.601) (1.891,0.617) (1.922,0.634) (1.954,0.650) (1.986,0.667) (2.017,0.684) (2.048,0.701) (2.080,0.718) (2.111,0.735) (2.142,0.753) (2.173,0.770) (2.204,0.788) (2.235,0.806) (2.266,0.823) (2.296,0.841) (2.327,0.860) (2.358,0.878) (2.388,0.896) (2.419,0.915) (2.449,0.934) (2.479,0.952) (2.510,0.971) (2.540,0.990) (2.570,1.009) (2.600,1.029)};
\addplot[lineaE, flechaen=0.289] coordinates {(-2.600,1.029) (-2.570,1.009) (-2.540,0.990) (-2.510,0.971) (-2.479,0.952) (-2.449,0.934) (-2.419,0.915) (-2.388,0.896) (-2.358,0.878) (-2.327,0.860) (-2.296,0.841) (-2.266,0.823) (-2.235,0.806) (-2.204,0.788) (-2.173,0.770) (-2.142,0.753) (-2.111,0.735) (-2.080,0.718) (-2.048,0.701) (-2.017,0.684) (-1.986,0.667) (-1.954,0.650) (-1.922,0.634) (-1.891,0.617) (-1.859,0.601) (-1.827,0.585) (-1.795,0.569) (-1.764,0.553) (-1.732,0.537) (-1.699,0.521) (-1.667,0.506) (-1.635,0.491) (-1.603,0.476) (-1.570,0.461) (-1.538,0.446) (-1.505,0.431) (-1.473,0.417) (-1.440,0.402) (-1.407,0.388) (-1.375,0.374) (-1.342,0.360) (-1.309,0.347) (-1.276,0.333) (-1.243,0.320) (-1.210,0.307) (-1.176,0.294) (-1.143,0.281) (-1.110,0.268) (-1.076,0.256) (-1.043,0.244) (-1.009,0.232) (-0.975,0.221)};
\addplot[lineaE, flechaen=0.874] coordinates {(0.925,0.381) (0.954,0.400) (0.984,0.420) (1.012,0.441) (1.041,0.462) (1.070,0.483) (1.098,0.505) (1.126,0.527) (1.154,0.549) (1.181,0.571) (1.209,0.594) (1.236,0.617) (1.263,0.640) (1.289,0.663) (1.316,0.687) (1.342,0.711) (1.368,0.735) (1.394,0.760) (1.419,0.785) (1.444,0.810) (1.469,0.835) (1.494,0.860) (1.519,0.886) (1.543,0.912) (1.567,0.938) (1.591,0.964) (1.615,0.991) (1.638,1.018) (1.661,1.044) (1.684,1.072) (1.707,1.099) (1.729,1.127) (1.751,1.154) (1.773,1.182) (1.795,1.210) (1.816,1.239) (1.838,1.267) (1.858,1.296) (1.879,1.325) (1.900,1.354) (1.920,1.383) (1.940,1.412) (1.959,1.442) (1.979,1.472) (1.998,1.502) (2.017,1.532) (2.036,1.562) (2.054,1.592) (2.072,1.623) (2.090,1.653) (2.108,1.684) (2.125,1.715) (2.143,1.746) (2.159,1.778) (2.176,1.809) (2.193,1.840) (2.209,1.872) (2.224,1.904) (2.240,1.936) (2.255,1.968) (2.270,2.000)};
\addplot[lineaE, flechaen=0.126] coordinates {(-2.270,2.000) (-2.255,1.968) (-2.240,1.936) (-2.224,1.904) (-2.209,1.872) (-2.193,1.840) (-2.176,1.809) (-2.159,1.778) (-2.143,1.746) (-2.125,1.715) (-2.108,1.684) (-2.090,1.653) (-2.072,1.623) (-2.054,1.592) (-2.036,1.562) (-2.017,1.532) (-1.998,1.502) (-1.979,1.472) (-1.959,1.442) (-1.940,1.412) (-1.920,1.383) (-1.900,1.354) (-1.879,1.325) (-1.858,1.296) (-1.838,1.267) (-1.816,1.239) (-1.795,1.210) (-1.773,1.182) (-1.751,1.154) (-1.729,1.127) (-1.707,1.099) (-1.684,1.072) (-1.661,1.044) (-1.638,1.018) (-1.615,0.991) (-1.591,0.964) (-1.567,0.938) (-1.543,0.912) (-1.519,0.886) (-1.494,0.860) (-1.469,0.835) (-1.444,0.810) (-1.419,0.785) (-1.394,0.760) (-1.368,0.735) (-1.342,0.711) (-1.316,0.687) (-1.289,0.663) (-1.263,0.640) (-1.236,0.617) (-1.209,0.594) (-1.181,0.571) (-1.154,0.549) (-1.126,0.527) (-1.098,0.505) (-1.070,0.483) (-1.041,0.462) (-1.012,0.441) (-0.984,0.420) (-0.954,0.400) (-0.925,0.381)};
\addplot[lineaE, flechaen=0.500] coordinates {(0.854,0.520) (0.877,0.547) (0.900,0.573) (0.923,0.601) (0.944,0.628) (0.966,0.656) (0.987,0.684) (1.007,0.713) (1.028,0.742) (1.047,0.771) (1.066,0.801) (1.085,0.831) (1.103,0.861) (1.121,0.891) (1.138,0.922) (1.155,0.953) (1.171,0.984) (1.187,1.015) (1.203,1.047) (1.218,1.079) (1.232,1.111) (1.246,1.143) (1.259,1.176) (1.272,1.209) (1.284,1.242) (1.296,1.275) (1.308,1.308) (1.319,1.342) (1.329,1.375) (1.339,1.409) (1.348,1.443) (1.357,1.477) (1.365,1.511) (1.373,1.546) (1.380,1.580) (1.387,1.615) (1.393,1.650) (1.398,1.684) (1.403,1.719) (1.408,1.754) (1.412,1.789) (1.415,1.824) (1.418,1.859) (1.420,1.894) (1.422,1.930) (1.423,1.965) (1.423,2.000)};
\addplot[lineaE, flechaen=0.500] coordinates {(-1.423,2.000) (-1.423,1.965) (-1.422,1.930) (-1.420,1.894) (-1.418,1.859) (-1.415,1.824) (-1.412,1.789) (-1.408,1.754) (-1.403,1.719) (-1.398,1.684) (-1.393,1.650) (-1.387,1.615) (-1.380,1.580) (-1.373,1.546) (-1.365,1.511) (-1.357,1.477) (-1.348,1.443) (-1.339,1.409) (-1.329,1.375) (-1.319,1.342) (-1.308,1.308) (-1.296,1.275) (-1.284,1.242) (-1.272,1.209) (-1.259,1.176) (-1.246,1.143) (-1.232,1.111) (-1.218,1.079) (-1.203,1.047) (-1.187,1.015) (-1.171,0.984) (-1.155,0.953) (-1.138,0.922) (-1.121,0.891) (-1.103,0.861) (-1.085,0.831) (-1.066,0.801) (-1.047,0.771) (-1.028,0.742) (-1.007,0.713) (-0.987,0.684) (-0.966,0.656) (-0.944,0.628) (-0.923,0.601) (-0.900,0.573) (-0.877,0.547) (-0.854,0.520)};
\addplot[lineaE, flechaen=0.500] coordinates {(0.768,0.641) (0.785,0.672) (0.800,0.704) (0.815,0.736) (0.828,0.769) (0.841,0.802) (0.854,0.836) (0.865,0.869) (0.876,0.903) (0.886,0.937) (0.895,0.972) (0.903,1.006) (0.911,1.041) (0.917,1.076) (0.923,1.111) (0.928,1.146) (0.932,1.181) (0.935,1.217) (0.938,1.252) (0.939,1.288) (0.940,1.323) (0.939,1.359) (0.938,1.395) (0.936,1.430) (0.933,1.465) (0.929,1.501) (0.924,1.536) (0.918,1.571) (0.911,1.606) (0.904,1.641) (0.895,1.675) (0.885,1.709) (0.875,1.743) (0.863,1.777) (0.850,1.810) (0.837,1.843) (0.822,1.875) (0.807,1.907) (0.790,1.939) (0.773,1.970) (0.754,2.000)};
\addplot[lineaE, flechaen=0.500] coordinates {(-0.754,2.000) (-0.773,1.970) (-0.790,1.939) (-0.807,1.907) (-0.822,1.875) (-0.837,1.843) (-0.850,1.810) (-0.863,1.777) (-0.875,1.743) (-0.885,1.709) (-0.895,1.675) (-0.904,1.641) (-0.911,1.606) (-0.918,1.571) (-0.924,1.536) (-0.929,1.501) (-0.933,1.465) (-0.936,1.430) (-0.938,1.395) (-0.939,1.359) (-0.940,1.323) (-0.939,1.288) (-0.938,1.252) (-0.935,1.217) (-0.932,1.181) (-0.928,1.146) (-0.923,1.111) (-0.917,1.076) (-0.911,1.041) (-0.903,1.006) (-0.895,0.972) (-0.886,0.937) (-0.876,0.903) (-0.865,0.869) (-0.854,0.836) (-0.841,0.802) (-0.828,0.769) (-0.815,0.736) (-0.800,0.704) (-0.785,0.672) (-0.768,0.641)};
\addplot[lineaE, flechaen=0.707] coordinates {(0.661,0.751) (0.668,0.785) (0.674,0.820) (0.678,0.855) (0.681,0.890) (0.683,0.926) (0.684,0.961) (0.684,0.996) (0.682,1.032) (0.679,1.067) (0.675,1.102) (0.670,1.137) (0.663,1.172) (0.655,1.206) (0.646,1.240) (0.635,1.274) (0.623,1.307) (0.610,1.340) (0.596,1.372) (0.580,1.404) (0.562,1.435) (0.544,1.465) (0.524,1.494) (0.503,1.522) (0.480,1.550) (0.456,1.576) (0.431,1.601) (0.405,1.624) (0.377,1.646) (0.349,1.667) (0.319,1.686) (0.288,1.704) (0.257,1.720) (0.224,1.734) (0.191,1.746) (0.157,1.757) (0.123,1.765) (0.088,1.771) (0.053,1.775) (0.018,1.778) (-0.018,1.778) (-0.053,1.775) (-0.088,1.771) (-0.123,1.765) (-0.157,1.757) (-0.191,1.746) (-0.224,1.734) (-0.257,1.720) (-0.288,1.704) (-0.319,1.686) (-0.349,1.667) (-0.377,1.646) (-0.405,1.624) (-0.431,1.601) (-0.456,1.576) (-0.480,1.550) (-0.503,1.522) (-0.524,1.494) (-0.544,1.465) (-0.562,1.435) (-0.580,1.404) (-0.596,1.372) (-0.610,1.340) (-0.623,1.307) (-0.635,1.274) (-0.646,1.240) (-0.655,1.206) (-0.663,1.172) (-0.670,1.137) (-0.675,1.102) (-0.679,1.067) (-0.682,1.032) (-0.684,0.996) (-0.684,0.961) (-0.683,0.926) (-0.681,0.890) (-0.678,0.855) (-0.674,0.820) (-0.668,0.785) (-0.661,0.751)};
\addplot[lineaE, flechaen=0.077] coordinates {(0.525,0.851) (0.521,0.886) (0.515,0.921) (0.506,0.955) (0.495,0.989) (0.483,1.023) (0.470,1.056) (0.454,1.088) (0.437,1.119) (0.418,1.149) (0.397,1.178) (0.374,1.206) (0.350,1.232) (0.324,1.256) (0.297,1.279) (0.268,1.300) (0.238,1.319) (0.206,1.336) (0.174,1.350) (0.140,1.362) (0.106,1.372) (0.071,1.379) (0.036,1.383) (-0.000,1.384) (-0.036,1.383) (-0.071,1.379) (-0.106,1.372) (-0.140,1.362) (-0.174,1.350) (-0.206,1.336) (-0.238,1.319) (-0.268,1.300) (-0.297,1.279) (-0.324,1.256) (-0.350,1.232) (-0.374,1.206) (-0.397,1.178) (-0.418,1.149) (-0.437,1.119) (-0.454,1.088) (-0.470,1.056) (-0.483,1.023) (-0.495,0.989) (-0.506,0.955) (-0.515,0.921) (-0.521,0.886) (-0.525,0.851)};
\addplot[lineaE, flechaen=0.500] coordinates {(0.340,0.941) (0.320,0.969) (0.297,0.996) (0.272,1.021) (0.245,1.044) (0.217,1.065) (0.186,1.083) (0.155,1.099) (0.122,1.112) (0.088,1.121) (0.053,1.128) (0.018,1.131) (-0.018,1.131) (-0.053,1.128) (-0.088,1.121) (-0.122,1.112) (-0.155,1.099) (-0.186,1.083) (-0.217,1.065) (-0.245,1.044) (-0.272,1.021) (-0.297,0.996) (-0.320,0.969) (-0.340,0.941)};
\addplot[lineaE, flechaen=0.500] coordinates {(0.335,0.942) (0.301,0.940) (0.266,0.940) (0.230,0.940) (0.195,0.940) (0.159,0.940) (0.124,0.940) (0.089,0.940) (0.053,0.940) (0.018,0.940) (-0.018,0.940) (-0.053,0.940) (-0.089,0.940) (-0.124,0.940) (-0.159,0.940) (-0.195,0.940) (-0.230,0.940) (-0.266,0.940) (-0.301,0.940) (-0.335,0.942)};
\addplot[lineaE, flechaen=0.500] coordinates {(0.522,0.853) (0.488,0.850) (0.451,0.850) (0.415,0.850) (0.379,0.850) (0.343,0.850) (0.307,0.850) (0.271,0.850) (0.235,0.850) (0.199,0.850) (0.163,0.850) (0.126,0.850) (0.090,0.850) (0.054,0.850) (0.018,0.850) (-0.018,0.850) (-0.054,0.850) (-0.090,0.850) (-0.126,0.850) (-0.163,0.850) (-0.199,0.850) (-0.235,0.850) (-0.271,0.850) (-0.307,0.850) (-0.343,0.850) (-0.379,0.850) (-0.415,0.850) (-0.451,0.850) (-0.488,0.850) (-0.522,0.853)};
\addplot[lineaE, flechaen=0.121] coordinates {(0.660,0.751) (0.625,0.750) (0.589,0.750) (0.553,0.750) (0.518,0.750) (0.482,0.750) (0.446,0.750) (0.411,0.750) (0.375,0.750) (0.339,0.750) (0.303,0.750) (0.268,0.750) (0.232,0.750) (0.196,0.750) (0.161,0.750) (0.125,0.750) (0.089,0.750) (0.054,0.750) (0.018,0.750) (-0.018,0.750) (-0.054,0.750) (-0.089,0.750) (-0.125,0.750) (-0.161,0.750) (-0.196,0.750) (-0.232,0.750) (-0.268,0.750) (-0.303,0.750) (-0.339,0.750) (-0.375,0.750) (-0.411,0.750) (-0.446,0.750) (-0.482,0.750) (-0.518,0.750) (-0.553,0.750) (-0.589,0.750) (-0.625,0.750) (-0.660,0.751)};
\addplot[lineaE, flechaen=0.826] coordinates {(0.768,0.641) (0.732,0.640) (0.697,0.640) (0.661,0.640) (0.625,0.640) (0.589,0.640) (0.554,0.640) (0.518,0.640) (0.482,0.640) (0.447,0.640) (0.411,0.640) (0.375,0.640) (0.339,0.640) (0.304,0.640) (0.268,0.640) (0.232,0.640) (0.196,0.640) (0.161,0.640) (0.125,0.640) (0.089,0.640) (0.054,0.640) (0.018,0.640) (-0.018,0.640) (-0.054,0.640) (-0.089,0.640) (-0.125,0.640) (-0.161,0.640) (-0.196,0.640) (-0.232,0.640) (-0.268,0.640) (-0.304,0.640) (-0.339,0.640) (-0.375,0.640) (-0.411,0.640) (-0.447,0.640) (-0.482,0.640) (-0.518,0.640) (-0.554,0.640) (-0.589,0.640) (-0.625,0.640) (-0.661,0.640) (-0.697,0.640) (-0.732,0.640) (-0.768,0.641)};
\addplot[lineaE, flechaen=0.207] coordinates {(0.853,0.521) (0.818,0.520) (0.783,0.520) (0.747,0.520) (0.711,0.520) (0.676,0.520) (0.640,0.520) (0.605,0.520) (0.569,0.520) (0.534,0.520) (0.498,0.520) (0.462,0.520) (0.427,0.520) (0.391,0.520) (0.356,0.520) (0.320,0.520) (0.285,0.520) (0.249,0.520) (0.213,0.520) (0.178,0.520) (0.142,0.520) (0.107,0.520) (0.071,0.520) (0.036,0.520) (0.000,0.520) (-0.036,0.520) (-0.071,0.520) (-0.107,0.520) (-0.142,0.520) (-0.178,0.520) (-0.213,0.520) (-0.249,0.520) (-0.285,0.520) (-0.320,0.520) (-0.356,0.520) (-0.391,0.520) (-0.427,0.520) (-0.462,0.520) (-0.498,0.520) (-0.534,0.520) (-0.569,0.520) (-0.605,0.520) (-0.640,0.520) (-0.676,0.520) (-0.711,0.520) (-0.747,0.520) (-0.783,0.520) (-0.818,0.520) (-0.853,0.521)};
\addplot[lineaE, flechaen=0.770] coordinates {(0.925,0.381) (0.889,0.380) (0.854,0.380) (0.818,0.380) (0.783,0.380) (0.747,0.380) (0.711,0.380) (0.676,0.380) (0.640,0.380) (0.605,0.380) (0.569,0.380) (0.534,0.380) (0.498,0.380) (0.462,0.380) (0.427,0.380) (0.391,0.380) (0.356,0.380) (0.320,0.380) (0.285,0.380) (0.249,0.380) (0.213,0.380) (0.178,0.380) (0.142,0.380) (0.107,0.380) (0.071,0.380) (0.036,0.380) (-0.000,0.380) (-0.036,0.380) (-0.071,0.380) (-0.107,0.380) (-0.142,0.380) (-0.178,0.380) (-0.213,0.380) (-0.249,0.380) (-0.285,0.380) (-0.320,0.380) (-0.356,0.380) (-0.391,0.380) (-0.427,0.380) (-0.462,0.380) (-0.498,0.380) (-0.534,0.380) (-0.569,0.380) (-0.605,0.380) (-0.640,0.380) (-0.676,0.380) (-0.711,0.380) (-0.747,0.380) (-0.783,0.380) (-0.818,0.380) (-0.854,0.380) (-0.889,0.380) (-0.925,0.381)};
\addplot[lineaE, flechaen=0.244] coordinates {(0.975,0.220) (0.940,0.220) (0.905,0.220) (0.869,0.220) (0.834,0.220) (0.798,0.220) (0.763,0.220) (0.727,0.220) (0.692,0.220) (0.656,0.220) (0.621,0.220) (0.585,0.220) (0.550,0.220) (0.514,0.220) (0.479,0.220) (0.443,0.220) (0.408,0.220) (0.372,0.220) (0.337,0.220) (0.302,0.220) (0.266,0.220) (0.231,0.220) (0.195,0.220) (0.160,0.220) (0.124,0.220) (0.089,0.220) (0.053,0.220) (0.018,0.220) (-0.018,0.220) (-0.053,0.220) (-0.089,0.220) (-0.124,0.220) (-0.160,0.220) (-0.195,0.220) (-0.231,0.220) (-0.266,0.220) (-0.302,0.220) (-0.337,0.220) (-0.372,0.220) (-0.408,0.220) (-0.443,0.220) (-0.479,0.220) (-0.514,0.220) (-0.550,0.220) (-0.585,0.220) (-0.621,0.220) (-0.656,0.220) (-0.692,0.220) (-0.727,0.220) (-0.763,0.220) (-0.798,0.220) (-0.834,0.220) (-0.869,0.220) (-0.905,0.220) (-0.940,0.220) (-0.975,0.220)};
\addplot[lineaE, flechaen=0.711] coordinates {(0.975,-0.221) (1.009,-0.232) (1.043,-0.244) (1.076,-0.256) (1.110,-0.268) (1.143,-0.281) (1.176,-0.294) (1.210,-0.307) (1.243,-0.320) (1.276,-0.333) (1.309,-0.347) (1.342,-0.360) (1.375,-0.374) (1.407,-0.388) (1.440,-0.402) (1.473,-0.417) (1.505,-0.431) (1.538,-0.446) (1.570,-0.461) (1.603,-0.476) (1.635,-0.491) (1.667,-0.506) (1.699,-0.521) (1.732,-0.537) (1.764,-0.553) (1.795,-0.569) (1.827,-0.585) (1.859,-0.601) (1.891,-0.617) (1.922,-0.634) (1.954,-0.650) (1.986,-0.667) (2.017,-0.684) (2.048,-0.701) (2.080,-0.718) (2.111,-0.735) (2.142,-0.753) (2.173,-0.770) (2.204,-0.788) (2.235,-0.806) (2.266,-0.823) (2.296,-0.841) (2.327,-0.860) (2.358,-0.878) (2.388,-0.896) (2.419,-0.915) (2.449,-0.934) (2.479,-0.952) (2.510,-0.971) (2.540,-0.990) (2.570,-1.009) (2.600,-1.029)};
\addplot[lineaE, flechaen=0.289] coordinates {(-2.600,-1.029) (-2.570,-1.009) (-2.540,-0.990) (-2.510,-0.971) (-2.479,-0.952) (-2.449,-0.934) (-2.419,-0.915) (-2.388,-0.896) (-2.358,-0.878) (-2.327,-0.860) (-2.296,-0.841) (-2.266,-0.823) (-2.235,-0.806) (-2.204,-0.788) (-2.173,-0.770) (-2.142,-0.753) (-2.111,-0.735) (-2.080,-0.718) (-2.048,-0.701) (-2.017,-0.684) (-1.986,-0.667) (-1.954,-0.650) (-1.922,-0.634) (-1.891,-0.617) (-1.859,-0.601) (-1.827,-0.585) (-1.795,-0.569) (-1.764,-0.553) (-1.732,-0.537) (-1.699,-0.521) (-1.667,-0.506) (-1.635,-0.491) (-1.603,-0.476) (-1.570,-0.461) (-1.538,-0.446) (-1.505,-0.431) (-1.473,-0.417) (-1.440,-0.402) (-1.407,-0.388) (-1.375,-0.374) (-1.342,-0.360) (-1.309,-0.347) (-1.276,-0.333) (-1.243,-0.320) (-1.210,-0.307) (-1.176,-0.294) (-1.143,-0.281) (-1.110,-0.268) (-1.076,-0.256) (-1.043,-0.244) (-1.009,-0.232) (-0.975,-0.221)};
\addplot[lineaE, flechaen=0.874] coordinates {(0.925,-0.381) (0.954,-0.400) (0.984,-0.420) (1.012,-0.441) (1.041,-0.462) (1.070,-0.483) (1.098,-0.505) (1.126,-0.527) (1.154,-0.549) (1.181,-0.571) (1.209,-0.594) (1.236,-0.617) (1.263,-0.640) (1.289,-0.663) (1.316,-0.687) (1.342,-0.711) (1.368,-0.735) (1.394,-0.760) (1.419,-0.785) (1.444,-0.810) (1.469,-0.835) (1.494,-0.860) (1.519,-0.886) (1.543,-0.912) (1.567,-0.938) (1.591,-0.964) (1.615,-0.991) (1.638,-1.018) (1.661,-1.044) (1.684,-1.072) (1.707,-1.099) (1.729,-1.127) (1.751,-1.154) (1.773,-1.182) (1.795,-1.210) (1.816,-1.239) (1.838,-1.267) (1.858,-1.296) (1.879,-1.325) (1.900,-1.354) (1.920,-1.383) (1.940,-1.412) (1.959,-1.442) (1.979,-1.472) (1.998,-1.502) (2.017,-1.532) (2.036,-1.562) (2.054,-1.592) (2.072,-1.623) (2.090,-1.653) (2.108,-1.684) (2.125,-1.715) (2.143,-1.746) (2.159,-1.778) (2.176,-1.809) (2.193,-1.840) (2.209,-1.872) (2.224,-1.904) (2.240,-1.936) (2.255,-1.968) (2.270,-2.000)};
\addplot[lineaE, flechaen=0.126] coordinates {(-2.270,-2.000) (-2.255,-1.968) (-2.240,-1.936) (-2.224,-1.904) (-2.209,-1.872) (-2.193,-1.840) (-2.176,-1.809) (-2.159,-1.778) (-2.143,-1.746) (-2.125,-1.715) (-2.108,-1.684) (-2.090,-1.653) (-2.072,-1.623) (-2.054,-1.592) (-2.036,-1.562) (-2.017,-1.532) (-1.998,-1.502) (-1.979,-1.472) (-1.959,-1.442) (-1.940,-1.412) (-1.920,-1.383) (-1.900,-1.354) (-1.879,-1.325) (-1.858,-1.296) (-1.838,-1.267) (-1.816,-1.239) (-1.795,-1.210) (-1.773,-1.182) (-1.751,-1.154) (-1.729,-1.127) (-1.707,-1.099) (-1.684,-1.072) (-1.661,-1.044) (-1.638,-1.018) (-1.615,-0.991) (-1.591,-0.964) (-1.567,-0.938) (-1.543,-0.912) (-1.519,-0.886) (-1.494,-0.860) (-1.469,-0.835) (-1.444,-0.810) (-1.419,-0.785) (-1.394,-0.760) (-1.368,-0.735) (-1.342,-0.711) (-1.316,-0.687) (-1.289,-0.663) (-1.263,-0.640) (-1.236,-0.617) (-1.209,-0.594) (-1.181,-0.571) (-1.154,-0.549) (-1.126,-0.527) (-1.098,-0.505) (-1.070,-0.483) (-1.041,-0.462) (-1.012,-0.441) (-0.984,-0.420) (-0.954,-0.400) (-0.925,-0.381)};
\addplot[lineaE, flechaen=0.500] coordinates {(0.854,-0.520) (0.877,-0.547) (0.900,-0.573) (0.923,-0.601) (0.944,-0.628) (0.966,-0.656) (0.987,-0.684) (1.007,-0.713) (1.028,-0.742) (1.047,-0.771) (1.066,-0.801) (1.085,-0.831) (1.103,-0.861) (1.121,-0.891) (1.138,-0.922) (1.155,-0.953) (1.171,-0.984) (1.187,-1.015) (1.203,-1.047) (1.218,-1.079) (1.232,-1.111) (1.246,-1.143) (1.259,-1.176) (1.272,-1.209) (1.284,-1.242) (1.296,-1.275) (1.308,-1.308) (1.319,-1.342) (1.329,-1.375) (1.339,-1.409) (1.348,-1.443) (1.357,-1.477) (1.365,-1.511) (1.373,-1.546) (1.380,-1.580) (1.387,-1.615) (1.393,-1.650) (1.398,-1.684) (1.403,-1.719) (1.408,-1.754) (1.412,-1.789) (1.415,-1.824) (1.418,-1.859) (1.420,-1.894) (1.422,-1.930) (1.423,-1.965) (1.423,-2.000)};
\addplot[lineaE, flechaen=0.500] coordinates {(-1.423,-2.000) (-1.423,-1.965) (-1.422,-1.930) (-1.420,-1.894) (-1.418,-1.859) (-1.415,-1.824) (-1.412,-1.789) (-1.408,-1.754) (-1.403,-1.719) (-1.398,-1.684) (-1.393,-1.650) (-1.387,-1.615) (-1.380,-1.580) (-1.373,-1.546) (-1.365,-1.511) (-1.357,-1.477) (-1.348,-1.443) (-1.339,-1.409) (-1.329,-1.375) (-1.319,-1.342) (-1.308,-1.308) (-1.296,-1.275) (-1.284,-1.242) (-1.272,-1.209) (-1.259,-1.176) (-1.246,-1.143) (-1.232,-1.111) (-1.218,-1.079) (-1.203,-1.047) (-1.187,-1.015) (-1.171,-0.984) (-1.155,-0.953) (-1.138,-0.922) (-1.121,-0.891) (-1.103,-0.861) (-1.085,-0.831) (-1.066,-0.801) (-1.047,-0.771) (-1.028,-0.742) (-1.007,-0.713) (-0.987,-0.684) (-0.966,-0.656) (-0.944,-0.628) (-0.923,-0.601) (-0.900,-0.573) (-0.877,-0.547) (-0.854,-0.520)};
\addplot[lineaE, flechaen=0.500] coordinates {(0.768,-0.641) (0.785,-0.672) (0.800,-0.704) (0.815,-0.736) (0.828,-0.769) (0.841,-0.802) (0.854,-0.836) (0.865,-0.869) (0.876,-0.903) (0.886,-0.937) (0.895,-0.972) (0.903,-1.006) (0.911,-1.041) (0.917,-1.076) (0.923,-1.111) (0.928,-1.146) (0.932,-1.181) (0.935,-1.217) (0.938,-1.252) (0.939,-1.288) (0.940,-1.323) (0.939,-1.359) (0.938,-1.395) (0.936,-1.430) (0.933,-1.465) (0.929,-1.501) (0.924,-1.536) (0.918,-1.571) (0.911,-1.606) (0.904,-1.641) (0.895,-1.675) (0.885,-1.709) (0.875,-1.743) (0.863,-1.777) (0.850,-1.810) (0.837,-1.843) (0.822,-1.875) (0.807,-1.907) (0.790,-1.939) (0.773,-1.970) (0.754,-2.000)};
\addplot[lineaE, flechaen=0.500] coordinates {(-0.754,-2.000) (-0.773,-1.970) (-0.790,-1.939) (-0.807,-1.907) (-0.822,-1.875) (-0.837,-1.843) (-0.850,-1.810) (-0.863,-1.777) (-0.875,-1.743) (-0.885,-1.709) (-0.895,-1.675) (-0.904,-1.641) (-0.911,-1.606) (-0.918,-1.571) (-0.924,-1.536) (-0.929,-1.501) (-0.933,-1.465) (-0.936,-1.430) (-0.938,-1.395) (-0.939,-1.359) (-0.940,-1.323) (-0.939,-1.288) (-0.938,-1.252) (-0.935,-1.217) (-0.932,-1.181) (-0.928,-1.146) (-0.923,-1.111) (-0.917,-1.076) (-0.911,-1.041) (-0.903,-1.006) (-0.895,-0.972) (-0.886,-0.937) (-0.876,-0.903) (-0.865,-0.869) (-0.854,-0.836) (-0.841,-0.802) (-0.828,-0.769) (-0.815,-0.736) (-0.800,-0.704) (-0.785,-0.672) (-0.768,-0.641)};
\addplot[lineaE, flechaen=0.293] coordinates {(0.661,-0.751) (0.668,-0.785) (0.674,-0.820) (0.678,-0.855) (0.681,-0.890) (0.683,-0.926) (0.684,-0.961) (0.684,-0.996) (0.682,-1.032) (0.679,-1.067) (0.675,-1.102) (0.670,-1.137) (0.663,-1.172) (0.655,-1.206) (0.646,-1.240) (0.635,-1.274) (0.623,-1.307) (0.610,-1.340) (0.596,-1.372) (0.580,-1.404) (0.562,-1.435) (0.544,-1.465) (0.524,-1.494) (0.503,-1.522) (0.480,-1.550) (0.456,-1.576) (0.431,-1.601) (0.405,-1.624) (0.377,-1.646) (0.349,-1.667) (0.319,-1.686) (0.288,-1.704) (0.257,-1.720) (0.224,-1.734) (0.191,-1.746) (0.157,-1.757) (0.123,-1.765) (0.088,-1.771) (0.053,-1.775) (0.018,-1.778) (-0.018,-1.778) (-0.053,-1.775) (-0.088,-1.771) (-0.123,-1.765) (-0.157,-1.757) (-0.191,-1.746) (-0.224,-1.734) (-0.257,-1.720) (-0.288,-1.704) (-0.319,-1.686) (-0.349,-1.667) (-0.377,-1.646) (-0.405,-1.624) (-0.431,-1.601) (-0.456,-1.576) (-0.480,-1.550) (-0.503,-1.522) (-0.524,-1.494) (-0.544,-1.465) (-0.562,-1.435) (-0.580,-1.404) (-0.596,-1.372) (-0.610,-1.340) (-0.623,-1.307) (-0.635,-1.274) (-0.646,-1.240) (-0.655,-1.206) (-0.663,-1.172) (-0.670,-1.137) (-0.675,-1.102) (-0.679,-1.067) (-0.682,-1.032) (-0.684,-0.996) (-0.684,-0.961) (-0.683,-0.926) (-0.681,-0.890) (-0.678,-0.855) (-0.674,-0.820) (-0.668,-0.785) (-0.661,-0.751)};
\addplot[lineaE, flechaen=0.923] coordinates {(0.525,-0.851) (0.521,-0.886) (0.515,-0.921) (0.506,-0.955) (0.495,-0.989) (0.483,-1.023) (0.470,-1.056) (0.454,-1.088) (0.437,-1.119) (0.418,-1.149) (0.397,-1.178) (0.374,-1.206) (0.350,-1.232) (0.324,-1.256) (0.297,-1.279) (0.268,-1.300) (0.238,-1.319) (0.206,-1.336) (0.174,-1.350) (0.140,-1.362) (0.106,-1.372) (0.071,-1.379) (0.036,-1.383) (-0.000,-1.384) (-0.036,-1.383) (-0.071,-1.379) (-0.106,-1.372) (-0.140,-1.362) (-0.174,-1.350) (-0.206,-1.336) (-0.238,-1.319) (-0.268,-1.300) (-0.297,-1.279) (-0.324,-1.256) (-0.350,-1.232) (-0.374,-1.206) (-0.397,-1.178) (-0.418,-1.149) (-0.437,-1.119) (-0.454,-1.088) (-0.470,-1.056) (-0.483,-1.023) (-0.495,-0.989) (-0.506,-0.955) (-0.515,-0.921) (-0.521,-0.886) (-0.525,-0.851)};
\addplot[lineaE, flechaen=0.500] coordinates {(0.340,-0.941) (0.320,-0.969) (0.297,-0.996) (0.272,-1.021) (0.245,-1.044) (0.217,-1.065) (0.186,-1.083) (0.155,-1.099) (0.122,-1.112) (0.088,-1.121) (0.053,-1.128) (0.018,-1.131) (-0.018,-1.131) (-0.053,-1.128) (-0.088,-1.121) (-0.122,-1.112) (-0.155,-1.099) (-0.186,-1.083) (-0.217,-1.065) (-0.245,-1.044) (-0.272,-1.021) (-0.297,-0.996) (-0.320,-0.969) (-0.340,-0.941)};
\addplot[lineaE, flechaen=0.500] coordinates {(0.335,-0.942) (0.301,-0.940) (0.266,-0.940) (0.230,-0.940) (0.195,-0.940) (0.159,-0.940) (0.124,-0.940) (0.089,-0.940) (0.053,-0.940) (0.018,-0.940) (-0.018,-0.940) (-0.053,-0.940) (-0.089,-0.940) (-0.124,-0.940) (-0.159,-0.940) (-0.195,-0.940) (-0.230,-0.940) (-0.266,-0.940) (-0.301,-0.940) (-0.335,-0.942)};
\addplot[lineaE, flechaen=0.500] coordinates {(0.522,-0.853) (0.488,-0.850) (0.451,-0.850) (0.415,-0.850) (0.379,-0.850) (0.343,-0.850) (0.307,-0.850) (0.271,-0.850) (0.235,-0.850) (0.199,-0.850) (0.163,-0.850) (0.126,-0.850) (0.090,-0.850) (0.054,-0.850) (0.018,-0.850) (-0.018,-0.850) (-0.054,-0.850) (-0.090,-0.850) (-0.126,-0.850) (-0.163,-0.850) (-0.199,-0.850) (-0.235,-0.850) (-0.271,-0.850) (-0.307,-0.850) (-0.343,-0.850) (-0.379,-0.850) (-0.415,-0.850) (-0.451,-0.850) (-0.488,-0.850) (-0.522,-0.853)};
\addplot[lineaE, flechaen=0.879] coordinates {(0.660,-0.751) (0.625,-0.750) (0.589,-0.750) (0.553,-0.750) (0.518,-0.750) (0.482,-0.750) (0.446,-0.750) (0.411,-0.750) (0.375,-0.750) (0.339,-0.750) (0.303,-0.750) (0.268,-0.750) (0.232,-0.750) (0.196,-0.750) (0.161,-0.750) (0.125,-0.750) (0.089,-0.750) (0.054,-0.750) (0.018,-0.750) (-0.018,-0.750) (-0.054,-0.750) (-0.089,-0.750) (-0.125,-0.750) (-0.161,-0.750) (-0.196,-0.750) (-0.232,-0.750) (-0.268,-0.750) (-0.303,-0.750) (-0.339,-0.750) (-0.375,-0.750) (-0.411,-0.750) (-0.446,-0.750) (-0.482,-0.750) (-0.518,-0.750) (-0.553,-0.750) (-0.589,-0.750) (-0.625,-0.750) (-0.660,-0.751)};
\addplot[lineaE, flechaen=0.174] coordinates {(0.768,-0.641) (0.732,-0.640) (0.697,-0.640) (0.661,-0.640) (0.625,-0.640) (0.589,-0.640) (0.554,-0.640) (0.518,-0.640) (0.482,-0.640) (0.447,-0.640) (0.411,-0.640) (0.375,-0.640) (0.339,-0.640) (0.304,-0.640) (0.268,-0.640) (0.232,-0.640) (0.196,-0.640) (0.161,-0.640) (0.125,-0.640) (0.089,-0.640) (0.054,-0.640) (0.018,-0.640) (-0.018,-0.640) (-0.054,-0.640) (-0.089,-0.640) (-0.125,-0.640) (-0.161,-0.640) (-0.196,-0.640) (-0.232,-0.640) (-0.268,-0.640) (-0.304,-0.640) (-0.339,-0.640) (-0.375,-0.640) (-0.411,-0.640) (-0.447,-0.640) (-0.482,-0.640) (-0.518,-0.640) (-0.554,-0.640) (-0.589,-0.640) (-0.625,-0.640) (-0.661,-0.640) (-0.697,-0.640) (-0.732,-0.640) (-0.768,-0.641)};
\addplot[lineaE, flechaen=0.793] coordinates {(0.853,-0.521) (0.818,-0.520) (0.783,-0.520) (0.747,-0.520) (0.711,-0.520) (0.676,-0.520) (0.640,-0.520) (0.605,-0.520) (0.569,-0.520) (0.534,-0.520) (0.498,-0.520) (0.462,-0.520) (0.427,-0.520) (0.391,-0.520) (0.356,-0.520) (0.320,-0.520) (0.285,-0.520) (0.249,-0.520) (0.213,-0.520) (0.178,-0.520) (0.142,-0.520) (0.107,-0.520) (0.071,-0.520) (0.036,-0.520) (0.000,-0.520) (-0.036,-0.520) (-0.071,-0.520) (-0.107,-0.520) (-0.142,-0.520) (-0.178,-0.520) (-0.213,-0.520) (-0.249,-0.520) (-0.285,-0.520) (-0.320,-0.520) (-0.356,-0.520) (-0.391,-0.520) (-0.427,-0.520) (-0.462,-0.520) (-0.498,-0.520) (-0.534,-0.520) (-0.569,-0.520) (-0.605,-0.520) (-0.640,-0.520) (-0.676,-0.520) (-0.711,-0.520) (-0.747,-0.520) (-0.783,-0.520) (-0.818,-0.520) (-0.853,-0.521)};
\addplot[lineaE, flechaen=0.230] coordinates {(0.925,-0.381) (0.889,-0.380) (0.854,-0.380) (0.818,-0.380) (0.783,-0.380) (0.747,-0.380) (0.711,-0.380) (0.676,-0.380) (0.640,-0.380) (0.605,-0.380) (0.569,-0.380) (0.534,-0.380) (0.498,-0.380) (0.462,-0.380) (0.427,-0.380) (0.391,-0.380) (0.356,-0.380) (0.320,-0.380) (0.285,-0.380) (0.249,-0.380) (0.213,-0.380) (0.178,-0.380) (0.142,-0.380) (0.107,-0.380) (0.071,-0.380) (0.036,-0.380) (-0.000,-0.380) (-0.036,-0.380) (-0.071,-0.380) (-0.107,-0.380) (-0.142,-0.380) (-0.178,-0.380) (-0.213,-0.380) (-0.249,-0.380) (-0.285,-0.380) (-0.320,-0.380) (-0.356,-0.380) (-0.391,-0.380) (-0.427,-0.380) (-0.462,-0.380) (-0.498,-0.380) (-0.534,-0.380) (-0.569,-0.380) (-0.605,-0.380) (-0.640,-0.380) (-0.676,-0.380) (-0.711,-0.380) (-0.747,-0.380) (-0.783,-0.380) (-0.818,-0.380) (-0.854,-0.380) (-0.889,-0.380) (-0.925,-0.381)};
\addplot[lineaE, flechaen=0.756] coordinates {(0.975,-0.220) (0.940,-0.220) (0.905,-0.220) (0.869,-0.220) (0.834,-0.220) (0.798,-0.220) (0.763,-0.220) (0.727,-0.220) (0.692,-0.220) (0.656,-0.220) (0.621,-0.220) (0.585,-0.220) (0.550,-0.220) (0.514,-0.220) (0.479,-0.220) (0.443,-0.220) (0.408,-0.220) (0.372,-0.220) (0.337,-0.220) (0.302,-0.220) (0.266,-0.220) (0.231,-0.220) (0.195,-0.220) (0.160,-0.220) (0.124,-0.220) (0.089,-0.220) (0.053,-0.220) (0.018,-0.220) (-0.018,-0.220) (-0.053,-0.220) (-0.089,-0.220) (-0.124,-0.220) (-0.160,-0.220) (-0.195,-0.220) (-0.231,-0.220) (-0.266,-0.220) (-0.302,-0.220) (-0.337,-0.220) (-0.372,-0.220) (-0.408,-0.220) (-0.443,-0.220) (-0.479,-0.220) (-0.514,-0.220) (-0.550,-0.220) (-0.585,-0.220) (-0.621,-0.220) (-0.656,-0.220) (-0.692,-0.220) (-0.727,-0.220) (-0.763,-0.220) (-0.798,-0.220) (-0.834,-0.220) (-0.869,-0.220) (-0.905,-0.220) (-0.940,-0.220) (-0.975,-0.220)};
\addplot[lineaE, flechaen=0.281] coordinates {(-2.600,0.000) (-1.000,0.000)};
\addplot[lineaE, flechaen=0.719] coordinates {(1.000,0.000) (2.600,0.000)};
\addplot[lineaE, flechaen=0.750] coordinates {(1.000,0.000) (-1.000,0.000)};

  \draw[texto, line width=0.6pt] (axis cs:0,0) circle[radius=1];
  \pgfplotsinvokeforeach{-60,-30,0,30,60}{
    \node[cpos, font=\scriptsize\bfseries, fill=fondo, circle, inner sep=0.2pt] at (axis cs:{1.16*cos(#1)},{1.16*sin(#1)}) {$+$};
    \node[cneg, font=\scriptsize\bfseries, fill=fondo, circle, inner sep=0.2pt] at (axis cs:{-1.16*cos(#1)},{1.16*sin(#1)}) {$-$};
  }
\end{axis}
\end{tikzpicture}
\caption{Líneas de $\vE$}
\end{subfigure}\hfill
\begin{subfigure}[b]{0.49\textwidth}\centering
\begin{tikzpicture}
\begin{axis}[esferap]
  \fill[dielec] (axis cs:0,0) circle[radius=1];
  %% Generado por scripts/lineas_campo.py -- no editar a mano
\addplot[equipot] coordinates {(-0.905,-0.429) (-0.926,-0.401) (-0.946,-0.372) (-0.964,-0.342) (-0.981,-0.311) (-0.996,-0.280) (-1.009,-0.247) (-1.021,-0.214) (-1.031,-0.180) (-1.039,-0.146) (-1.045,-0.112) (-1.050,-0.077) (-1.053,-0.042) (-1.054,-0.007) (-1.054,0.028) (-1.051,0.063) (-1.047,0.098) (-1.041,0.133) (-1.034,0.167) (-1.025,0.201) (-1.014,0.234) (-1.001,0.267) (-0.987,0.299) (-0.971,0.330) (-0.953,0.361) (-0.934,0.390) (-0.914,0.419) (-0.900,0.411) (-0.900,0.375) (-0.900,0.340) (-0.900,0.305) (-0.900,0.270) (-0.900,0.235) (-0.900,0.200) (-0.900,0.165) (-0.900,0.130) (-0.900,0.095) (-0.900,0.059) (-0.900,0.024) (-0.900,-0.011) (-0.900,-0.046) (-0.900,-0.081) (-0.900,-0.116) (-0.900,-0.151) (-0.900,-0.186) (-0.900,-0.221) (-0.900,-0.257) (-0.900,-0.292) (-0.900,-0.327) (-0.900,-0.362) (-0.900,-0.397) (-0.905,-0.429)};
\addplot[equipot] coordinates {(-0.726,-0.692) (-0.759,-0.678) (-0.790,-0.663) (-0.821,-0.646) (-0.852,-0.627) (-0.881,-0.607) (-0.908,-0.585) (-0.935,-0.562) (-0.961,-0.538) (-0.985,-0.512) (-1.008,-0.485) (-1.030,-0.457) (-1.050,-0.428) (-1.069,-0.398) (-1.087,-0.367) (-1.103,-0.336) (-1.117,-0.304) (-1.130,-0.271) (-1.142,-0.237) (-1.152,-0.203) (-1.160,-0.169) (-1.167,-0.134) (-1.172,-0.099) (-1.176,-0.064) (-1.178,-0.029) (-1.178,0.007) (-1.177,0.042) (-1.175,0.077) (-1.170,0.112) (-1.165,0.147) (-1.157,0.182) (-1.148,0.216) (-1.138,0.250) (-1.126,0.283) (-1.112,0.316) (-1.097,0.348) (-1.080,0.379) (-1.062,0.409) (-1.043,0.439) (-1.022,0.468) (-1.000,0.495) (-0.976,0.522) (-0.951,0.547) (-0.925,0.571) (-0.898,0.593) (-0.870,0.615) (-0.840,0.634) (-0.810,0.652) (-0.779,0.669) (-0.746,0.683) (-0.720,0.683) (-0.720,0.648) (-0.720,0.613) (-0.720,0.577) (-0.720,0.542) (-0.720,0.506) (-0.720,0.471) (-0.720,0.436) (-0.720,0.400) (-0.720,0.365) (-0.720,0.330) (-0.720,0.294) (-0.720,0.259) (-0.720,0.223) (-0.720,0.188) (-0.720,0.153) (-0.720,0.117) (-0.720,0.082) (-0.720,0.046) (-0.720,0.011) (-0.720,-0.024) (-0.720,-0.060) (-0.720,-0.095) (-0.720,-0.130) (-0.720,-0.166) (-0.720,-0.201) (-0.720,-0.237) (-0.720,-0.272) (-0.720,-0.307) (-0.720,-0.343) (-0.720,-0.378) (-0.720,-0.414) (-0.720,-0.449) (-0.720,-0.484) (-0.720,-0.520) (-0.720,-0.555) (-0.720,-0.590) (-0.720,-0.626) (-0.720,-0.661) (-0.726,-0.692)};
\addplot[equipot] coordinates {(-0.662,-0.841) (-0.697,-0.836) (-0.732,-0.830) (-0.766,-0.822) (-0.799,-0.812) (-0.833,-0.801) (-0.865,-0.788) (-0.897,-0.773) (-0.928,-0.757) (-0.959,-0.740) (-0.988,-0.721) (-1.017,-0.700) (-1.045,-0.679) (-1.072,-0.656) (-1.097,-0.633) (-1.122,-0.608) (-1.146,-0.582) (-1.168,-0.555) (-1.190,-0.527) (-1.210,-0.498) (-1.229,-0.469) (-1.247,-0.439) (-1.263,-0.408) (-1.279,-0.376) (-1.293,-0.344) (-1.306,-0.311) (-1.317,-0.278) (-1.327,-0.244) (-1.336,-0.210) (-1.344,-0.176) (-1.350,-0.142) (-1.355,-0.107) (-1.358,-0.072) (-1.360,-0.037) (-1.361,-0.002) (-1.360,0.033) (-1.358,0.068) (-1.355,0.103) (-1.350,0.138) (-1.344,0.173) (-1.337,0.207) (-1.328,0.241) (-1.318,0.275) (-1.307,0.308) (-1.294,0.341) (-1.280,0.373) (-1.265,0.404) (-1.249,0.435) (-1.231,0.466) (-1.212,0.495) (-1.192,0.524) (-1.171,0.552) (-1.148,0.579) (-1.125,0.605) (-1.100,0.630) (-1.074,0.654) (-1.048,0.677) (-1.020,0.698) (-0.991,0.719) (-0.962,0.738) (-0.932,0.755) (-0.900,0.772) (-0.869,0.786) (-0.836,0.800) (-0.803,0.811) (-0.769,0.821) (-0.735,0.829) (-0.701,0.836) (-0.666,0.841) (-0.631,0.843) (-0.596,0.844) (-0.561,0.843) (-0.540,0.827) (-0.540,0.792) (-0.540,0.756) (-0.540,0.721) (-0.540,0.686) (-0.540,0.651) (-0.540,0.616) (-0.540,0.581) (-0.540,0.546) (-0.540,0.511) (-0.540,0.476) (-0.540,0.441) (-0.540,0.405) (-0.540,0.370) (-0.540,0.335) (-0.540,0.300) (-0.540,0.265) (-0.540,0.230) (-0.540,0.195) (-0.540,0.160) (-0.540,0.125) (-0.540,0.090) (-0.540,0.054) (-0.540,0.019) (-0.540,-0.016) (-0.540,-0.051) (-0.540,-0.086) (-0.540,-0.121) (-0.540,-0.156) (-0.540,-0.191) (-0.540,-0.226) (-0.540,-0.261) (-0.540,-0.297) (-0.540,-0.332) (-0.540,-0.367) (-0.540,-0.402) (-0.540,-0.437) (-0.540,-0.472) (-0.540,-0.507) (-0.540,-0.542) (-0.540,-0.577) (-0.540,-0.613) (-0.540,-0.648) (-0.540,-0.683) (-0.540,-0.718) (-0.540,-0.753) (-0.540,-0.788) (-0.540,-0.823) (-0.557,-0.843) (-0.592,-0.844) (-0.627,-0.844) (-0.662,-0.841)};
\addplot[equipot] coordinates {(-0.732,-1.034) (-0.767,-1.033) (-0.802,-1.031) (-0.837,-1.027) (-0.872,-1.021) (-0.906,-1.015) (-0.940,-1.006) (-0.974,-0.996) (-1.007,-0.985) (-1.040,-0.973) (-1.073,-0.959) (-1.104,-0.944) (-1.136,-0.928) (-1.166,-0.911) (-1.196,-0.892) (-1.225,-0.873) (-1.254,-0.852) (-1.281,-0.830) (-1.308,-0.808) (-1.334,-0.784) (-1.360,-0.759) (-1.384,-0.734) (-1.407,-0.708) (-1.430,-0.681) (-1.451,-0.653) (-1.472,-0.625) (-1.492,-0.595) (-1.510,-0.566) (-1.528,-0.535) (-1.545,-0.504) (-1.560,-0.473) (-1.575,-0.441) (-1.589,-0.408) (-1.601,-0.375) (-1.612,-0.342) (-1.623,-0.309) (-1.632,-0.275) (-1.640,-0.241) (-1.647,-0.206) (-1.653,-0.171) (-1.658,-0.137) (-1.662,-0.102) (-1.665,-0.067) (-1.666,-0.031) (-1.667,0.004) (-1.666,0.039) (-1.664,0.074) (-1.661,0.109) (-1.657,0.144) (-1.652,0.179) (-1.646,0.213) (-1.639,0.248) (-1.630,0.282) (-1.621,0.316) (-1.610,0.349) (-1.599,0.382) (-1.586,0.415) (-1.572,0.448) (-1.557,0.479) (-1.541,0.511) (-1.524,0.542) (-1.507,0.572) (-1.488,0.602) (-1.468,0.631) (-1.447,0.659) (-1.425,0.687) (-1.402,0.713) (-1.379,0.739) (-1.354,0.765) (-1.329,0.789) (-1.303,0.812) (-1.276,0.835) (-1.248,0.856) (-1.219,0.877) (-1.190,0.896) (-1.160,0.914) (-1.129,0.932) (-1.098,0.947) (-1.066,0.962) (-1.033,0.976) (-1.000,0.988) (-0.967,0.999) (-0.933,1.008) (-0.899,1.016) (-0.864,1.023) (-0.830,1.028) (-0.795,1.031) (-0.759,1.033) (-0.724,1.034) (-0.689,1.033) (-0.654,1.030) (-0.619,1.026) (-0.585,1.020) (-0.550,1.012) (-0.516,1.003) (-0.483,0.992) (-0.450,0.979) (-0.418,0.964) (-0.387,0.948) (-0.360,0.928) (-0.360,0.893) (-0.360,0.858) (-0.360,0.823) (-0.360,0.787) (-0.360,0.752) (-0.360,0.717) (-0.360,0.682) (-0.360,0.647) (-0.360,0.612) (-0.360,0.576) (-0.360,0.541) (-0.360,0.506) (-0.360,0.471) (-0.360,0.436) (-0.360,0.401) (-0.360,0.366) (-0.360,0.330) (-0.360,0.295) (-0.360,0.260) (-0.360,0.225) (-0.360,0.190) (-0.360,0.155) (-0.360,0.119) (-0.360,0.084) (-0.360,0.049) (-0.360,0.014) (-0.360,-0.021) (-0.360,-0.056) (-0.360,-0.092) (-0.360,-0.127) (-0.360,-0.162) (-0.360,-0.197) (-0.360,-0.232) (-0.360,-0.267) (-0.360,-0.303) (-0.360,-0.338) (-0.360,-0.373) (-0.360,-0.408) (-0.360,-0.443) (-0.360,-0.478) (-0.360,-0.514) (-0.360,-0.549) (-0.360,-0.584) (-0.360,-0.619) (-0.360,-0.654) (-0.360,-0.689) (-0.360,-0.725) (-0.360,-0.760) (-0.360,-0.795) (-0.360,-0.830) (-0.360,-0.865) (-0.360,-0.900) (-0.363,-0.935) (-0.393,-0.952) (-0.425,-0.968) (-0.457,-0.982) (-0.490,-0.994) (-0.523,-1.005) (-0.557,-1.014) (-0.592,-1.021) (-0.627,-1.027) (-0.661,-1.031) (-0.697,-1.033) (-0.732,-1.034)};
\addplot[equipot] coordinates {(-1.148,-1.456) (-1.183,-1.452) (-1.218,-1.447) (-1.252,-1.441) (-1.287,-1.434) (-1.321,-1.425) (-1.355,-1.416) (-1.388,-1.406) (-1.421,-1.395) (-1.454,-1.383) (-1.487,-1.370) (-1.519,-1.356) (-1.551,-1.341) (-1.582,-1.325) (-1.613,-1.309) (-1.644,-1.291) (-1.674,-1.273) (-1.703,-1.254) (-1.732,-1.234) (-1.761,-1.214) (-1.789,-1.193) (-1.816,-1.171) (-1.843,-1.148) (-1.869,-1.125) (-1.895,-1.101) (-1.920,-1.077) (-1.945,-1.051) (-1.968,-1.026) (-1.992,-0.999) (-2.014,-0.973) (-2.036,-0.945) (-2.057,-0.917) (-2.078,-0.889) (-2.098,-0.860) (-2.117,-0.831) (-2.136,-0.801) (-2.153,-0.771) (-2.170,-0.740) (-2.187,-0.709) (-2.203,-0.677) (-2.217,-0.646) (-2.232,-0.614) (-2.245,-0.581) (-2.258,-0.549) (-2.270,-0.516) (-2.281,-0.482) (-2.291,-0.449) (-2.301,-0.415) (-2.310,-0.381) (-2.318,-0.347) (-2.326,-0.313) (-2.332,-0.278) (-2.338,-0.244) (-2.343,-0.209) (-2.347,-0.174) (-2.351,-0.139) (-2.354,-0.104) (-2.356,-0.069) (-2.357,-0.034) (-2.357,0.001) (-2.357,0.036) (-2.355,0.071) (-2.353,0.106) (-2.351,0.141) (-2.347,0.176) (-2.343,0.211) (-2.338,0.246) (-2.332,0.280) (-2.325,0.315) (-2.318,0.349) (-2.310,0.383) (-2.301,0.417) (-2.291,0.451) (-2.280,0.484) (-2.269,0.517) (-2.257,0.550) (-2.244,0.583) (-2.231,0.615) (-2.217,0.647) (-2.202,0.679) (-2.186,0.711) (-2.170,0.742) (-2.152,0.772) (-2.135,0.802) (-2.116,0.832) (-2.097,0.862) (-2.077,0.890) (-2.056,0.919) (-2.035,0.947) (-2.013,0.974) (-1.990,1.001) (-1.967,1.027) (-1.943,1.053) (-1.919,1.078) (-1.894,1.102) (-1.868,1.126) (-1.842,1.150) (-1.815,1.172) (-1.787,1.194) (-1.759,1.215) (-1.731,1.235) (-1.702,1.255) (-1.672,1.274) (-1.642,1.292) (-1.611,1.309) (-1.580,1.326) (-1.549,1.342) (-1.517,1.356) (-1.485,1.370) (-1.452,1.383) (-1.420,1.395) (-1.386,1.406) (-1.353,1.417) (-1.319,1.426) (-1.285,1.434) (-1.250,1.441) (-1.216,1.447) (-1.181,1.453) (-1.146,1.457) (-1.111,1.460) (-1.076,1.461) (-1.041,1.462) (-1.006,1.462) (-0.971,1.460) (-0.936,1.458) (-0.901,1.454) (-0.866,1.449) (-0.832,1.443) (-0.798,1.436) (-0.763,1.427) (-0.730,1.418) (-0.696,1.407) (-0.663,1.395) (-0.631,1.382) (-0.599,1.367) (-0.567,1.352) (-0.536,1.335) (-0.506,1.318) (-0.477,1.299) (-0.448,1.279) (-0.420,1.258) (-0.392,1.236) (-0.366,1.213) (-0.340,1.189) (-0.316,1.164) (-0.292,1.138) (-0.269,1.111) (-0.247,1.084) (-0.226,1.056) (-0.207,1.027) (-0.188,0.997) (-0.180,0.964) (-0.180,0.929) (-0.180,0.894) (-0.180,0.859) (-0.180,0.824) (-0.180,0.789) (-0.180,0.753) (-0.180,0.718) (-0.180,0.683) (-0.180,0.648) (-0.180,0.613) (-0.180,0.578) (-0.180,0.543) (-0.180,0.508) (-0.180,0.473) (-0.180,0.438) (-0.180,0.403) (-0.180,0.367) (-0.180,0.332) (-0.180,0.297) (-0.180,0.262) (-0.180,0.227) (-0.180,0.192) (-0.180,0.157) (-0.180,0.122) (-0.180,0.087) (-0.180,0.052) (-0.180,0.017) (-0.180,-0.019) (-0.180,-0.054) (-0.180,-0.089) (-0.180,-0.124) (-0.180,-0.159) (-0.180,-0.194) (-0.180,-0.229) (-0.180,-0.264) (-0.180,-0.299) (-0.180,-0.334) (-0.180,-0.369) (-0.180,-0.404) (-0.180,-0.440) (-0.180,-0.475) (-0.180,-0.510) (-0.180,-0.545) (-0.180,-0.580) (-0.180,-0.615) (-0.180,-0.650) (-0.180,-0.685) (-0.180,-0.720) (-0.180,-0.755) (-0.180,-0.790) (-0.180,-0.826) (-0.180,-0.861) (-0.180,-0.896) (-0.180,-0.931) (-0.180,-0.966) (-0.189,-0.999) (-0.208,-1.028) (-0.228,-1.057) (-0.248,-1.085) (-0.270,-1.113) (-0.293,-1.139) (-0.317,-1.165) (-0.342,-1.190) (-0.367,-1.214) (-0.394,-1.237) (-0.421,-1.259) (-0.449,-1.280) (-0.478,-1.300) (-0.508,-1.319) (-0.538,-1.336) (-0.569,-1.353) (-0.601,-1.368) (-0.633,-1.383) (-0.665,-1.396) (-0.698,-1.408) (-0.732,-1.418) (-0.765,-1.428) (-0.799,-1.436) (-0.834,-1.443) (-0.868,-1.449) (-0.903,-1.454) (-0.938,-1.458) (-0.973,-1.461) (-1.008,-1.462) (-1.043,-1.462) (-1.078,-1.461) (-1.113,-1.459) (-1.148,-1.456)};
\addplot[equipot] coordinates {(0.923,-1.457) (0.958,-1.460) (0.993,-1.462) (1.028,-1.462) (1.063,-1.462) (1.098,-1.460) (1.133,-1.458) (1.168,-1.454) (1.202,-1.450) (1.237,-1.444) (1.271,-1.437) (1.306,-1.429) (1.340,-1.420) (1.373,-1.411) (1.407,-1.400) (1.440,-1.388) (1.473,-1.375) (1.505,-1.362) (1.537,-1.347) (1.568,-1.332) (1.600,-1.316) (1.630,-1.299) (1.661,-1.281) (1.690,-1.263) (1.720,-1.243) (1.748,-1.223) (1.776,-1.202) (1.804,-1.181) (1.831,-1.158) (1.858,-1.135) (1.884,-1.112) (1.909,-1.087) (1.934,-1.063) (1.958,-1.037) (1.981,-1.011) (2.004,-0.984) (2.027,-0.957) (2.048,-0.930) (2.069,-0.901) (2.089,-0.873) (2.109,-0.844) (2.127,-0.814) (2.146,-0.784) (2.163,-0.753) (2.180,-0.723) (2.196,-0.691) (2.211,-0.660) (2.225,-0.628) (2.239,-0.596) (2.252,-0.563) (2.265,-0.530) (2.276,-0.497) (2.287,-0.464) (2.297,-0.430) (2.306,-0.396) (2.315,-0.362) (2.322,-0.328) (2.329,-0.293) (2.336,-0.259) (2.341,-0.224) (2.346,-0.189) (2.349,-0.155) (2.352,-0.120) (2.355,-0.085) (2.356,-0.050) (2.357,-0.014) (2.357,0.021) (2.356,0.056) (2.354,0.091) (2.352,0.126) (2.349,0.161) (2.345,0.196) (2.340,0.230) (2.335,0.265) (2.328,0.299) (2.321,0.334) (2.313,0.368) (2.305,0.402) (2.295,0.436) (2.285,0.469) (2.274,0.503) (2.262,0.536) (2.250,0.569) (2.237,0.601) (2.223,0.633) (2.208,0.665) (2.193,0.697) (2.177,0.728) (2.160,0.759) (2.143,0.789) (2.124,0.819) (2.105,0.849) (2.086,0.878) (2.065,0.906) (2.044,0.934) (2.023,0.962) (2.000,0.989) (1.977,1.016) (1.954,1.042) (1.930,1.067) (1.905,1.092) (1.879,1.116) (1.853,1.139) (1.827,1.162) (1.799,1.184) (1.772,1.206) (1.743,1.227) (1.714,1.247) (1.685,1.266) (1.655,1.284) (1.625,1.302) (1.594,1.319) (1.563,1.335) (1.531,1.350) (1.499,1.364) (1.467,1.378) (1.434,1.390) (1.401,1.402) (1.367,1.412) (1.334,1.422) (1.300,1.431) (1.266,1.438) (1.231,1.445) (1.196,1.450) (1.162,1.455) (1.127,1.458) (1.092,1.461) (1.057,1.462) (1.022,1.462) (0.986,1.461) (0.951,1.459) (0.917,1.456) (0.882,1.452) (0.847,1.446) (0.813,1.439) (0.778,1.431) (0.745,1.422) (0.711,1.412) (0.678,1.400) (0.645,1.388) (0.613,1.374) (0.581,1.359) (0.550,1.343) (0.519,1.326) (0.490,1.307) (0.460,1.288) (0.432,1.267) (0.404,1.246) (0.377,1.223) (0.351,1.199) (0.326,1.175) (0.302,1.149) (0.279,1.123) (0.257,1.096) (0.236,1.068) (0.215,1.039) (0.196,1.010) (0.180,0.979) (0.180,0.944) (0.180,0.909) (0.180,0.874) (0.180,0.839) (0.180,0.804) (0.180,0.769) (0.180,0.734) (0.180,0.699) (0.180,0.664) (0.180,0.629) (0.180,0.593) (0.180,0.558) (0.180,0.523) (0.180,0.488) (0.180,0.453) (0.180,0.418) (0.180,0.383) (0.180,0.348) (0.180,0.313) (0.180,0.278) (0.180,0.243) (0.180,0.207) (0.180,0.172) (0.180,0.137) (0.180,0.102) (0.180,0.067) (0.180,0.032) (0.180,-0.003) (0.180,-0.038) (0.180,-0.073) (0.180,-0.108) (0.180,-0.143) (0.180,-0.178) (0.180,-0.214) (0.180,-0.249) (0.180,-0.284) (0.180,-0.319) (0.180,-0.354) (0.180,-0.389) (0.180,-0.424) (0.180,-0.459) (0.180,-0.494) (0.180,-0.529) (0.180,-0.564) (0.180,-0.600) (0.180,-0.635) (0.180,-0.670) (0.180,-0.705) (0.180,-0.740) (0.180,-0.775) (0.180,-0.810) (0.180,-0.845) (0.180,-0.880) (0.180,-0.915) (0.180,-0.950) (0.181,-0.985) (0.199,-1.015) (0.219,-1.044) (0.239,-1.073) (0.261,-1.101) (0.283,-1.128) (0.306,-1.154) (0.331,-1.179) (0.356,-1.204) (0.382,-1.227) (0.409,-1.249) (0.437,-1.271) (0.465,-1.291) (0.495,-1.310) (0.525,-1.329) (0.555,-1.346) (0.587,-1.362) (0.618,-1.376) (0.651,-1.390) (0.684,-1.402) (0.717,-1.414) (0.750,-1.424) (0.784,-1.433) (0.819,-1.440) (0.853,-1.447) (0.888,-1.452) (0.923,-1.457)};
\addplot[equipot] coordinates {(0.732,-1.034) (0.767,-1.033) (0.802,-1.031) (0.837,-1.027) (0.872,-1.021) (0.906,-1.015) (0.940,-1.006) (0.974,-0.996) (1.007,-0.985) (1.040,-0.973) (1.073,-0.959) (1.104,-0.944) (1.136,-0.928) (1.166,-0.911) (1.196,-0.892) (1.225,-0.873) (1.254,-0.852) (1.281,-0.830) (1.308,-0.808) (1.334,-0.784) (1.360,-0.759) (1.384,-0.734) (1.407,-0.708) (1.430,-0.681) (1.451,-0.653) (1.472,-0.625) (1.492,-0.595) (1.510,-0.566) (1.528,-0.535) (1.545,-0.504) (1.560,-0.473) (1.575,-0.441) (1.589,-0.408) (1.601,-0.375) (1.612,-0.342) (1.623,-0.309) (1.632,-0.275) (1.640,-0.241) (1.647,-0.206) (1.653,-0.171) (1.658,-0.137) (1.662,-0.102) (1.665,-0.067) (1.666,-0.031) (1.667,0.004) (1.666,0.039) (1.664,0.074) (1.661,0.109) (1.657,0.144) (1.652,0.179) (1.646,0.213) (1.639,0.248) (1.630,0.282) (1.621,0.316) (1.610,0.349) (1.599,0.382) (1.586,0.415) (1.572,0.448) (1.557,0.479) (1.541,0.511) (1.524,0.542) (1.507,0.572) (1.488,0.602) (1.468,0.631) (1.447,0.659) (1.425,0.687) (1.402,0.713) (1.379,0.739) (1.354,0.765) (1.329,0.789) (1.303,0.812) (1.276,0.835) (1.248,0.856) (1.219,0.877) (1.190,0.896) (1.160,0.914) (1.129,0.932) (1.098,0.947) (1.066,0.962) (1.033,0.976) (1.000,0.988) (0.967,0.999) (0.933,1.008) (0.899,1.016) (0.864,1.023) (0.830,1.028) (0.795,1.031) (0.759,1.033) (0.724,1.034) (0.689,1.033) (0.654,1.030) (0.619,1.026) (0.585,1.020) (0.550,1.012) (0.516,1.003) (0.483,0.992) (0.450,0.979) (0.418,0.964) (0.387,0.948) (0.360,0.928) (0.360,0.893) (0.360,0.858) (0.360,0.823) (0.360,0.787) (0.360,0.752) (0.360,0.717) (0.360,0.682) (0.360,0.647) (0.360,0.612) (0.360,0.576) (0.360,0.541) (0.360,0.506) (0.360,0.471) (0.360,0.436) (0.360,0.401) (0.360,0.366) (0.360,0.330) (0.360,0.295) (0.360,0.260) (0.360,0.225) (0.360,0.190) (0.360,0.155) (0.360,0.119) (0.360,0.084) (0.360,0.049) (0.360,0.014) (0.360,-0.021) (0.360,-0.056) (0.360,-0.092) (0.360,-0.127) (0.360,-0.162) (0.360,-0.197) (0.360,-0.232) (0.360,-0.267) (0.360,-0.303) (0.360,-0.338) (0.360,-0.373) (0.360,-0.408) (0.360,-0.443) (0.360,-0.478) (0.360,-0.514) (0.360,-0.549) (0.360,-0.584) (0.360,-0.619) (0.360,-0.654) (0.360,-0.689) (0.360,-0.725) (0.360,-0.760) (0.360,-0.795) (0.360,-0.830) (0.360,-0.865) (0.360,-0.900) (0.363,-0.935) (0.393,-0.952) (0.425,-0.968) (0.457,-0.982) (0.490,-0.994) (0.523,-1.005) (0.557,-1.014) (0.592,-1.021) (0.627,-1.027) (0.661,-1.031) (0.697,-1.033) (0.732,-1.034)};
\addplot[equipot] coordinates {(0.541,-0.841) (0.576,-0.844) (0.611,-0.844) (0.646,-0.842) (0.681,-0.839) (0.716,-0.833) (0.750,-0.826) (0.784,-0.817) (0.817,-0.806) (0.850,-0.794) (0.882,-0.780) (0.914,-0.765) (0.945,-0.748) (0.975,-0.730) (1.004,-0.710) (1.032,-0.689) (1.059,-0.667) (1.086,-0.644) (1.111,-0.619) (1.135,-0.594) (1.158,-0.567) (1.180,-0.540) (1.201,-0.512) (1.220,-0.483) (1.239,-0.453) (1.256,-0.422) (1.272,-0.391) (1.286,-0.359) (1.300,-0.327) (1.312,-0.294) (1.323,-0.260) (1.332,-0.226) (1.340,-0.192) (1.347,-0.158) (1.352,-0.123) (1.357,-0.088) (1.359,-0.053) (1.361,-0.018) (1.361,0.017) (1.359,0.052) (1.357,0.087) (1.353,0.122) (1.347,0.157) (1.340,0.191) (1.332,0.225) (1.323,0.259) (1.312,0.292) (1.300,0.325) (1.287,0.358) (1.272,0.390) (1.256,0.421) (1.239,0.452) (1.221,0.482) (1.201,0.511) (1.181,0.539) (1.159,0.567) (1.136,0.593) (1.112,0.619) (1.086,0.643) (1.060,0.666) (1.033,0.688) (1.005,0.709) (0.976,0.729) (0.946,0.747) (0.915,0.764) (0.884,0.780) (0.851,0.794) (0.818,0.806) (0.785,0.817) (0.751,0.826) (0.717,0.833) (0.682,0.839) (0.647,0.842) (0.612,0.844) (0.577,0.844) (0.542,0.842) (0.540,0.808) (0.540,0.773) (0.540,0.738) (0.540,0.703) (0.540,0.668) (0.540,0.632) (0.540,0.597) (0.540,0.562) (0.540,0.527) (0.540,0.492) (0.540,0.457) (0.540,0.422) (0.540,0.387) (0.540,0.352) (0.540,0.317) (0.540,0.281) (0.540,0.246) (0.540,0.211) (0.540,0.176) (0.540,0.141) (0.540,0.106) (0.540,0.071) (0.540,0.036) (0.540,0.001) (0.540,-0.034) (0.540,-0.070) (0.540,-0.105) (0.540,-0.140) (0.540,-0.175) (0.540,-0.210) (0.540,-0.245) (0.540,-0.280) (0.540,-0.315) (0.540,-0.350) (0.540,-0.385) (0.540,-0.421) (0.540,-0.456) (0.540,-0.491) (0.540,-0.526) (0.540,-0.561) (0.540,-0.596) (0.540,-0.631) (0.540,-0.666) (0.540,-0.701) (0.540,-0.737) (0.540,-0.772) (0.540,-0.807) (0.541,-0.841)};
\addplot[equipot] coordinates {(0.720,-0.691) (0.753,-0.681) (0.785,-0.666) (0.816,-0.649) (0.847,-0.630) (0.876,-0.610) (0.904,-0.589) (0.931,-0.566) (0.957,-0.542) (0.981,-0.516) (1.004,-0.490) (1.026,-0.462) (1.047,-0.433) (1.066,-0.403) (1.084,-0.373) (1.100,-0.341) (1.115,-0.309) (1.128,-0.276) (1.140,-0.243) (1.150,-0.209) (1.159,-0.175) (1.166,-0.140) (1.171,-0.105) (1.175,-0.070) (1.178,-0.035) (1.179,0.001) (1.178,0.036) (1.175,0.071) (1.171,0.107) (1.166,0.141) (1.158,0.176) (1.150,0.210) (1.139,0.244) (1.128,0.278) (1.114,0.310) (1.099,0.342) (1.083,0.374) (1.065,0.404) (1.046,0.434) (1.025,0.463) (1.003,0.491) (0.980,0.517) (0.956,0.543) (0.930,0.567) (0.903,0.590) (0.875,0.611) (0.845,0.631) (0.815,0.649) (0.784,0.666) (0.752,0.681) (0.720,0.689) (0.720,0.654) (0.720,0.619) (0.720,0.583) (0.720,0.548) (0.720,0.512) (0.720,0.477) (0.720,0.442) (0.720,0.406) (0.720,0.371) (0.720,0.335) (0.720,0.300) (0.720,0.265) (0.720,0.229) (0.720,0.194) (0.720,0.159) (0.720,0.123) (0.720,0.088) (0.720,0.052) (0.720,0.017) (0.720,-0.018) (0.720,-0.054) (0.720,-0.089) (0.720,-0.125) (0.720,-0.160) (0.720,-0.195) (0.720,-0.231) (0.720,-0.266) (0.720,-0.301) (0.720,-0.337) (0.720,-0.372) (0.720,-0.408) (0.720,-0.443) (0.720,-0.478) (0.720,-0.514) (0.720,-0.549) (0.720,-0.585) (0.720,-0.620) (0.720,-0.655) (0.720,-0.691)};
\addplot[equipot] coordinates {(0.905,-0.429) (0.926,-0.401) (0.946,-0.372) (0.964,-0.342) (0.981,-0.311) (0.996,-0.280) (1.009,-0.247) (1.021,-0.214) (1.031,-0.180) (1.039,-0.146) (1.045,-0.112) (1.050,-0.077) (1.053,-0.042) (1.054,-0.007) (1.054,0.028) (1.051,0.063) (1.047,0.098) (1.041,0.133) (1.034,0.167) (1.025,0.201) (1.014,0.234) (1.001,0.267) (0.987,0.299) (0.971,0.330) (0.953,0.361) (0.934,0.390) (0.914,0.419) (0.900,0.411) (0.900,0.375) (0.900,0.340) (0.900,0.305) (0.900,0.270) (0.900,0.235) (0.900,0.200) (0.900,0.165) (0.900,0.130) (0.900,0.095) (0.900,0.059) (0.900,0.024) (0.900,-0.011) (0.900,-0.046) (0.900,-0.081) (0.900,-0.116) (0.900,-0.151) (0.900,-0.186) (0.900,-0.221) (0.900,-0.257) (0.900,-0.292) (0.900,-0.327) (0.900,-0.362) (0.900,-0.397) (0.905,-0.429)};
\addplot[lineaD, flechaen=0.094, flechaen=0.906] coordinates {(-2.600,1.029) (-2.570,1.010) (-2.541,0.991) (-2.511,0.972) (-2.481,0.953) (-2.451,0.935) (-2.421,0.916) (-2.391,0.898) (-2.361,0.880) (-2.331,0.862) (-2.301,0.844) (-2.271,0.826) (-2.240,0.809) (-2.210,0.791) (-2.179,0.774) (-2.149,0.756) (-2.118,0.739) (-2.087,0.722) (-2.056,0.705) (-2.026,0.688) (-1.995,0.672) (-1.964,0.655) (-1.932,0.639) (-1.901,0.623) (-1.870,0.606) (-1.839,0.590) (-1.807,0.575) (-1.776,0.559) (-1.744,0.543) (-1.713,0.528) (-1.681,0.513) (-1.649,0.498) (-1.618,0.483) (-1.586,0.468) (-1.554,0.453) (-1.522,0.438) (-1.490,0.424) (-1.458,0.410) (-1.425,0.396) (-1.393,0.382) (-1.361,0.368) (-1.328,0.355) (-1.296,0.341) (-1.263,0.328) (-1.231,0.315) (-1.198,0.302) (-1.165,0.289) (-1.132,0.277) (-1.099,0.265) (-1.066,0.252) (-1.033,0.240) (-1.000,0.229) (-0.967,0.220) (-0.931,0.220) (-0.896,0.220) (-0.861,0.220) (-0.826,0.220) (-0.791,0.220) (-0.756,0.220) (-0.721,0.220) (-0.685,0.220) (-0.650,0.220) (-0.615,0.220) (-0.580,0.220) (-0.545,0.220) (-0.510,0.220) (-0.475,0.220) (-0.439,0.220) (-0.404,0.220) (-0.369,0.220) (-0.334,0.220) (-0.299,0.220) (-0.264,0.220) (-0.228,0.220) (-0.193,0.220) (-0.158,0.220) (-0.123,0.220) (-0.088,0.220) (-0.053,0.220) (-0.018,0.220) (0.018,0.220) (0.053,0.220) (0.088,0.220) (0.123,0.220) (0.158,0.220) (0.193,0.220) (0.228,0.220) (0.264,0.220) (0.299,0.220) (0.334,0.220) (0.369,0.220) (0.404,0.220) (0.439,0.220) (0.475,0.220) (0.510,0.220) (0.545,0.220) (0.580,0.220) (0.615,0.220) (0.650,0.220) (0.685,0.220) (0.721,0.220) (0.756,0.220) (0.791,0.220) (0.826,0.220) (0.861,0.220) (0.896,0.220) (0.931,0.220) (0.967,0.220) (1.000,0.229) (1.033,0.240) (1.066,0.252) (1.099,0.265) (1.132,0.277) (1.165,0.289) (1.198,0.302) (1.231,0.315) (1.263,0.328) (1.296,0.341) (1.328,0.355) (1.361,0.368) (1.393,0.382) (1.425,0.396) (1.458,0.410) (1.490,0.424) (1.522,0.438) (1.554,0.453) (1.586,0.468) (1.618,0.483) (1.649,0.498) (1.681,0.513) (1.713,0.528) (1.744,0.543) (1.776,0.559) (1.807,0.575) (1.839,0.590) (1.870,0.606) (1.901,0.623) (1.932,0.639) (1.964,0.655) (1.995,0.672) (2.026,0.688) (2.056,0.705) (2.087,0.722) (2.118,0.739) (2.149,0.756) (2.179,0.774) (2.210,0.791) (2.240,0.809) (2.271,0.826) (2.301,0.844) (2.331,0.862) (2.361,0.880) (2.391,0.898) (2.421,0.916) (2.451,0.935) (2.481,0.953) (2.511,0.972) (2.541,0.991) (2.570,1.010) (2.600,1.029)};
\addplot[lineaD, flechaen=0.044, flechaen=0.956] coordinates {(-2.270,2.000) (-2.256,1.968) (-2.240,1.936) (-2.225,1.905) (-2.209,1.873) (-2.193,1.842) (-2.177,1.811) (-2.161,1.780) (-2.144,1.749) (-2.127,1.718) (-2.110,1.688) (-2.093,1.657) (-2.075,1.627) (-2.057,1.597) (-2.039,1.567) (-2.020,1.537) (-2.002,1.507) (-1.983,1.477) (-1.963,1.448) (-1.944,1.419) (-1.924,1.390) (-1.904,1.361) (-1.884,1.332) (-1.864,1.303) (-1.843,1.275) (-1.822,1.247) (-1.801,1.219) (-1.780,1.191) (-1.758,1.163) (-1.736,1.135) (-1.714,1.108) (-1.692,1.081) (-1.669,1.054) (-1.646,1.027) (-1.623,1.001) (-1.600,0.975) (-1.577,0.948) (-1.553,0.923) (-1.529,0.897) (-1.505,0.871) (-1.480,0.846) (-1.456,0.821) (-1.431,0.796) (-1.406,0.772) (-1.380,0.747) (-1.355,0.723) (-1.329,0.699) (-1.303,0.676) (-1.277,0.652) (-1.250,0.629) (-1.224,0.607) (-1.197,0.584) (-1.170,0.562) (-1.142,0.540) (-1.115,0.518) (-1.087,0.496) (-1.059,0.475) (-1.031,0.454) (-1.002,0.434) (-0.974,0.413) (-0.945,0.393) (-0.913,0.380) (-0.878,0.380) (-0.843,0.380) (-0.808,0.380) (-0.773,0.380) (-0.738,0.380) (-0.703,0.380) (-0.667,0.380) (-0.632,0.380) (-0.597,0.380) (-0.562,0.380) (-0.527,0.380) (-0.492,0.380) (-0.457,0.380) (-0.422,0.380) (-0.386,0.380) (-0.351,0.380) (-0.316,0.380) (-0.281,0.380) (-0.246,0.380) (-0.211,0.380) (-0.176,0.380) (-0.141,0.380) (-0.105,0.380) (-0.070,0.380) (-0.035,0.380) (-0.000,0.380) (0.035,0.380) (0.070,0.380) (0.105,0.380) (0.141,0.380) (0.176,0.380) (0.211,0.380) (0.246,0.380) (0.281,0.380) (0.316,0.380) (0.351,0.380) (0.386,0.380) (0.422,0.380) (0.457,0.380) (0.492,0.380) (0.527,0.380) (0.562,0.380) (0.597,0.380) (0.632,0.380) (0.667,0.380) (0.703,0.380) (0.738,0.380) (0.773,0.380) (0.808,0.380) (0.843,0.380) (0.878,0.380) (0.913,0.380) (0.945,0.393) (0.974,0.413) (1.002,0.434) (1.031,0.454) (1.059,0.475) (1.087,0.496) (1.115,0.518) (1.142,0.540) (1.170,0.562) (1.197,0.584) (1.224,0.607) (1.250,0.629) (1.277,0.652) (1.303,0.676) (1.329,0.699) (1.355,0.723) (1.380,0.747) (1.406,0.772) (1.431,0.796) (1.456,0.821) (1.480,0.846) (1.505,0.871) (1.529,0.897) (1.553,0.923) (1.577,0.948) (1.600,0.975) (1.623,1.001) (1.646,1.027) (1.669,1.054) (1.692,1.081) (1.714,1.108) (1.736,1.135) (1.758,1.163) (1.780,1.191) (1.801,1.219) (1.822,1.247) (1.843,1.275) (1.864,1.303) (1.884,1.332) (1.904,1.361) (1.924,1.390) (1.944,1.419) (1.963,1.448) (1.983,1.477) (2.002,1.507) (2.020,1.537) (2.039,1.567) (2.057,1.597) (2.075,1.627) (2.093,1.657) (2.110,1.688) (2.127,1.718) (2.144,1.749) (2.161,1.780) (2.177,1.811) (2.193,1.842) (2.209,1.873) (2.225,1.905) (2.240,1.936) (2.256,1.968) (2.270,2.000)};
\addplot[lineaD, flechaen=0.399] coordinates {(-1.423,2.000) (-1.423,1.965) (-1.422,1.930) (-1.420,1.895) (-1.418,1.860) (-1.415,1.825) (-1.412,1.790) (-1.408,1.755) (-1.404,1.720) (-1.399,1.685) (-1.393,1.651) (-1.387,1.616) (-1.380,1.582) (-1.373,1.547) (-1.366,1.513) (-1.357,1.479) (-1.349,1.445) (-1.339,1.411) (-1.330,1.378) (-1.319,1.344) (-1.309,1.311) (-1.297,1.277) (-1.285,1.244) (-1.273,1.212) (-1.260,1.179) (-1.247,1.146) (-1.233,1.114) (-1.219,1.082) (-1.204,1.050) (-1.189,1.019) (-1.173,0.987) (-1.157,0.956) (-1.140,0.925) (-1.123,0.895) (-1.105,0.864) (-1.087,0.834) (-1.069,0.805) (-1.050,0.775) (-1.030,0.746) (-1.010,0.717) (-0.990,0.689) (-0.969,0.660) (-0.948,0.632) (-0.926,0.605) (-0.904,0.578) (-0.881,0.551) (-0.858,0.524) (-0.825,0.520) (-0.789,0.520) (-0.754,0.520) (-0.719,0.520) (-0.684,0.520) (-0.649,0.520) (-0.614,0.520) (-0.579,0.520) (-0.544,0.520) (-0.509,0.520) (-0.474,0.520) (-0.439,0.520) (-0.404,0.520) (-0.368,0.520) (-0.333,0.520) (-0.298,0.520) (-0.263,0.520) (-0.228,0.520) (-0.193,0.520) (-0.158,0.520) (-0.123,0.520) (-0.088,0.520) (-0.053,0.520) (-0.018,0.520) (0.018,0.520) (0.053,0.520) (0.088,0.520) (0.123,0.520) (0.158,0.520) (0.193,0.520) (0.228,0.520) (0.263,0.520) (0.298,0.520) (0.333,0.520) (0.368,0.520) (0.404,0.520) (0.439,0.520) (0.474,0.520) (0.509,0.520) (0.544,0.520) (0.579,0.520) (0.614,0.520) (0.649,0.520) (0.684,0.520) (0.719,0.520) (0.754,0.520) (0.789,0.520) (0.825,0.520) (0.858,0.524) (0.881,0.551) (0.904,0.578) (0.926,0.605) (0.948,0.632) (0.969,0.660) (0.990,0.689) (1.010,0.717) (1.030,0.746) (1.050,0.775) (1.069,0.805) (1.087,0.834) (1.105,0.864) (1.123,0.895) (1.140,0.925) (1.157,0.956) (1.173,0.987) (1.189,1.019) (1.204,1.050) (1.219,1.082) (1.233,1.114) (1.247,1.146) (1.260,1.179) (1.273,1.212) (1.285,1.244) (1.297,1.277) (1.309,1.311) (1.319,1.344) (1.330,1.378) (1.339,1.411) (1.349,1.445) (1.357,1.479) (1.366,1.513) (1.373,1.547) (1.380,1.582) (1.387,1.616) (1.393,1.651) (1.399,1.685) (1.404,1.720) (1.408,1.755) (1.412,1.790) (1.415,1.825) (1.418,1.860) (1.420,1.895) (1.422,1.930) (1.423,1.965) (1.423,2.000)};
\addplot[lineaD, flechaen=0.614] coordinates {(-0.754,2.000) (-0.772,1.970) (-0.790,1.940) (-0.806,1.909) (-0.821,1.877) (-0.836,1.845) (-0.849,1.813) (-0.862,1.780) (-0.873,1.747) (-0.884,1.714) (-0.894,1.680) (-0.902,1.646) (-0.910,1.612) (-0.917,1.578) (-0.923,1.543) (-0.928,1.508) (-0.932,1.474) (-0.935,1.439) (-0.938,1.404) (-0.939,1.369) (-0.940,1.334) (-0.939,1.299) (-0.938,1.264) (-0.936,1.229) (-0.933,1.194) (-0.930,1.159) (-0.925,1.124) (-0.920,1.090) (-0.913,1.055) (-0.906,1.021) (-0.899,0.987) (-0.890,0.953) (-0.881,0.919) (-0.870,0.885) (-0.859,0.852) (-0.848,0.819) (-0.835,0.786) (-0.822,0.754) (-0.808,0.722) (-0.793,0.690) (-0.778,0.658) (-0.753,0.640) (-0.718,0.640) (-0.683,0.640) (-0.648,0.640) (-0.613,0.640) (-0.578,0.640) (-0.543,0.640) (-0.508,0.640) (-0.473,0.640) (-0.438,0.640) (-0.403,0.640) (-0.368,0.640) (-0.333,0.640) (-0.298,0.640) (-0.263,0.640) (-0.228,0.640) (-0.193,0.640) (-0.158,0.640) (-0.123,0.640) (-0.088,0.640) (-0.053,0.640) (-0.018,0.640) (0.018,0.640) (0.053,0.640) (0.088,0.640) (0.123,0.640) (0.158,0.640) (0.193,0.640) (0.228,0.640) (0.263,0.640) (0.298,0.640) (0.333,0.640) (0.368,0.640) (0.403,0.640) (0.438,0.640) (0.473,0.640) (0.508,0.640) (0.543,0.640) (0.578,0.640) (0.613,0.640) (0.648,0.640) (0.683,0.640) (0.718,0.640) (0.753,0.640) (0.778,0.658) (0.793,0.690) (0.808,0.722) (0.822,0.754) (0.835,0.786) (0.848,0.819) (0.859,0.852) (0.870,0.885) (0.881,0.919) (0.890,0.953) (0.899,0.987) (0.906,1.021) (0.913,1.055) (0.920,1.090) (0.925,1.124) (0.930,1.159) (0.933,1.194) (0.936,1.229) (0.938,1.264) (0.939,1.299) (0.940,1.334) (0.939,1.369) (0.938,1.404) (0.935,1.439) (0.932,1.474) (0.928,1.508) (0.923,1.543) (0.917,1.578) (0.910,1.612) (0.902,1.646) (0.894,1.680) (0.884,1.714) (0.873,1.747) (0.862,1.780) (0.849,1.813) (0.836,1.845) (0.821,1.877) (0.806,1.909) (0.790,1.940) (0.772,1.970) (0.754,2.000)};
\addplot[lineaD, flechaen=0.500] coordinates {(-0.657,0.750) (-0.621,0.750) (-0.586,0.750) (-0.551,0.750) (-0.516,0.750) (-0.481,0.750) (-0.446,0.750) (-0.410,0.750) (-0.375,0.750) (-0.340,0.750) (-0.305,0.750) (-0.270,0.750) (-0.234,0.750) (-0.199,0.750) (-0.164,0.750) (-0.129,0.750) (-0.094,0.750) (-0.059,0.750) (-0.023,0.750) (0.012,0.750) (0.047,0.750) (0.082,0.750) (0.117,0.750) (0.152,0.750) (0.188,0.750) (0.223,0.750) (0.258,0.750) (0.293,0.750) (0.328,0.750) (0.363,0.750) (0.399,0.750) (0.434,0.750) (0.469,0.750) (0.504,0.750) (0.539,0.750) (0.574,0.750) (0.610,0.750) (0.645,0.750) (0.665,0.770) (0.671,0.804) (0.676,0.839) (0.680,0.874) (0.683,0.909) (0.684,0.944) (0.684,0.979) (0.683,1.015) (0.681,1.050) (0.677,1.085) (0.673,1.119) (0.667,1.154) (0.660,1.189) (0.651,1.223) (0.641,1.256) (0.630,1.290) (0.617,1.323) (0.604,1.355) (0.589,1.387) (0.572,1.418) (0.554,1.448) (0.535,1.478) (0.515,1.506) (0.493,1.534) (0.470,1.561) (0.446,1.586) (0.420,1.610) (0.394,1.633) (0.366,1.655) (0.337,1.675) (0.307,1.693) (0.276,1.710) (0.245,1.725) (0.212,1.739) (0.179,1.750) (0.145,1.760) (0.111,1.767) (0.076,1.773) (0.041,1.776) (0.006,1.778) (-0.029,1.777) (-0.064,1.774) (-0.099,1.769) (-0.134,1.763) (-0.168,1.754) (-0.201,1.743) (-0.234,1.730) (-0.266,1.716) (-0.297,1.699) (-0.327,1.681) (-0.356,1.662) (-0.385,1.641) (-0.412,1.618) (-0.437,1.594) (-0.462,1.569) (-0.486,1.543) (-0.508,1.516) (-0.529,1.487) (-0.548,1.458) (-0.566,1.428) (-0.583,1.397) (-0.599,1.366) (-0.613,1.334) (-0.626,1.301) (-0.638,1.268) (-0.648,1.234) (-0.657,1.200) (-0.664,1.166) (-0.671,1.131) (-0.676,1.096) (-0.680,1.061) (-0.683,1.026) (-0.684,0.991) (-0.684,0.956) (-0.683,0.921) (-0.681,0.886) (-0.678,0.851) (-0.673,0.816) (-0.668,0.781) (-0.657,0.750)};
\addplot[lineaD, flechaen=0.382] coordinates {(-0.523,0.850) (-0.488,0.850) (-0.453,0.850) (-0.417,0.850) (-0.382,0.850) (-0.346,0.850) (-0.311,0.850) (-0.276,0.850) (-0.240,0.850) (-0.205,0.850) (-0.170,0.850) (-0.134,0.850) (-0.099,0.850) (-0.063,0.850) (-0.028,0.850) (0.007,0.850) (0.043,0.850) (0.078,0.850) (0.114,0.850) (0.149,0.850) (0.184,0.850) (0.220,0.850) (0.255,0.850) (0.291,0.850) (0.326,0.850) (0.361,0.850) (0.397,0.850) (0.432,0.850) (0.468,0.850) (0.503,0.850) (0.525,0.864) (0.519,0.898) (0.512,0.933) (0.502,0.967) (0.492,1.001) (0.479,1.034) (0.465,1.066) (0.449,1.098) (0.431,1.129) (0.411,1.158) (0.390,1.186) (0.367,1.213) (0.343,1.239) (0.317,1.263) (0.289,1.285) (0.260,1.305) (0.230,1.324) (0.198,1.340) (0.166,1.354) (0.132,1.365) (0.098,1.374) (0.063,1.380) (0.028,1.383) (-0.007,1.384) (-0.043,1.382) (-0.078,1.377) (-0.113,1.370) (-0.147,1.360) (-0.180,1.348) (-0.212,1.333) (-0.243,1.316) (-0.273,1.297) (-0.301,1.276) (-0.328,1.253) (-0.353,1.228) (-0.377,1.202) (-0.399,1.175) (-0.420,1.146) (-0.439,1.116) (-0.456,1.085) (-0.471,1.053) (-0.485,1.020) (-0.496,0.987) (-0.507,0.953) (-0.515,0.918) (-0.522,0.884) (-0.523,0.850)};
\addplot[lineaD, flechaen=0.500] coordinates {(-0.338,0.941) (-0.303,0.940) (-0.267,0.940) (-0.232,0.940) (-0.196,0.940) (-0.161,0.940) (-0.125,0.940) (-0.090,0.940) (-0.054,0.940) (-0.019,0.940) (0.017,0.940) (0.052,0.940) (0.088,0.940) (0.123,0.940) (0.159,0.940) (0.194,0.940) (0.230,0.940) (0.266,0.940) (0.301,0.940) (0.337,0.940) (0.321,0.967) (0.299,0.994) (0.274,1.020) (0.247,1.043) (0.218,1.064) (0.188,1.082) (0.156,1.098) (0.123,1.111) (0.089,1.121) (0.054,1.128) (0.019,1.131) (-0.017,1.131) (-0.052,1.128) (-0.087,1.122) (-0.121,1.112) (-0.154,1.099) (-0.186,1.083) (-0.217,1.065) (-0.245,1.044) (-0.272,1.021) (-0.297,0.996) (-0.320,0.969) (-0.338,0.941)};
\addplot[lineaD, flechaen=0.094, flechaen=0.906] coordinates {(-2.600,-1.029) (-2.570,-1.010) (-2.541,-0.991) (-2.511,-0.972) (-2.481,-0.953) (-2.451,-0.935) (-2.421,-0.916) (-2.391,-0.898) (-2.361,-0.880) (-2.331,-0.862) (-2.301,-0.844) (-2.271,-0.826) (-2.240,-0.809) (-2.210,-0.791) (-2.179,-0.774) (-2.149,-0.756) (-2.118,-0.739) (-2.087,-0.722) (-2.056,-0.705) (-2.026,-0.688) (-1.995,-0.672) (-1.964,-0.655) (-1.932,-0.639) (-1.901,-0.623) (-1.870,-0.606) (-1.839,-0.590) (-1.807,-0.575) (-1.776,-0.559) (-1.744,-0.543) (-1.713,-0.528) (-1.681,-0.513) (-1.649,-0.498) (-1.618,-0.483) (-1.586,-0.468) (-1.554,-0.453) (-1.522,-0.438) (-1.490,-0.424) (-1.458,-0.410) (-1.425,-0.396) (-1.393,-0.382) (-1.361,-0.368) (-1.328,-0.355) (-1.296,-0.341) (-1.263,-0.328) (-1.231,-0.315) (-1.198,-0.302) (-1.165,-0.289) (-1.132,-0.277) (-1.099,-0.265) (-1.066,-0.252) (-1.033,-0.240) (-1.000,-0.229) (-0.967,-0.220) (-0.931,-0.220) (-0.896,-0.220) (-0.861,-0.220) (-0.826,-0.220) (-0.791,-0.220) (-0.756,-0.220) (-0.721,-0.220) (-0.685,-0.220) (-0.650,-0.220) (-0.615,-0.220) (-0.580,-0.220) (-0.545,-0.220) (-0.510,-0.220) (-0.475,-0.220) (-0.439,-0.220) (-0.404,-0.220) (-0.369,-0.220) (-0.334,-0.220) (-0.299,-0.220) (-0.264,-0.220) (-0.228,-0.220) (-0.193,-0.220) (-0.158,-0.220) (-0.123,-0.220) (-0.088,-0.220) (-0.053,-0.220) (-0.018,-0.220) (0.018,-0.220) (0.053,-0.220) (0.088,-0.220) (0.123,-0.220) (0.158,-0.220) (0.193,-0.220) (0.228,-0.220) (0.264,-0.220) (0.299,-0.220) (0.334,-0.220) (0.369,-0.220) (0.404,-0.220) (0.439,-0.220) (0.475,-0.220) (0.510,-0.220) (0.545,-0.220) (0.580,-0.220) (0.615,-0.220) (0.650,-0.220) (0.685,-0.220) (0.721,-0.220) (0.756,-0.220) (0.791,-0.220) (0.826,-0.220) (0.861,-0.220) (0.896,-0.220) (0.931,-0.220) (0.967,-0.220) (1.000,-0.229) (1.033,-0.240) (1.066,-0.252) (1.099,-0.265) (1.132,-0.277) (1.165,-0.289) (1.198,-0.302) (1.231,-0.315) (1.263,-0.328) (1.296,-0.341) (1.328,-0.355) (1.361,-0.368) (1.393,-0.382) (1.425,-0.396) (1.458,-0.410) (1.490,-0.424) (1.522,-0.438) (1.554,-0.453) (1.586,-0.468) (1.618,-0.483) (1.649,-0.498) (1.681,-0.513) (1.713,-0.528) (1.744,-0.543) (1.776,-0.559) (1.807,-0.575) (1.839,-0.590) (1.870,-0.606) (1.901,-0.623) (1.932,-0.639) (1.964,-0.655) (1.995,-0.672) (2.026,-0.688) (2.056,-0.705) (2.087,-0.722) (2.118,-0.739) (2.149,-0.756) (2.179,-0.774) (2.210,-0.791) (2.240,-0.809) (2.271,-0.826) (2.301,-0.844) (2.331,-0.862) (2.361,-0.880) (2.391,-0.898) (2.421,-0.916) (2.451,-0.935) (2.481,-0.953) (2.511,-0.972) (2.541,-0.991) (2.570,-1.010) (2.600,-1.029)};
\addplot[lineaD, flechaen=0.044, flechaen=0.956] coordinates {(-2.270,-2.000) (-2.256,-1.968) (-2.240,-1.936) (-2.225,-1.905) (-2.209,-1.873) (-2.193,-1.842) (-2.177,-1.811) (-2.161,-1.780) (-2.144,-1.749) (-2.127,-1.718) (-2.110,-1.688) (-2.093,-1.657) (-2.075,-1.627) (-2.057,-1.597) (-2.039,-1.567) (-2.020,-1.537) (-2.002,-1.507) (-1.983,-1.477) (-1.963,-1.448) (-1.944,-1.419) (-1.924,-1.390) (-1.904,-1.361) (-1.884,-1.332) (-1.864,-1.303) (-1.843,-1.275) (-1.822,-1.247) (-1.801,-1.219) (-1.780,-1.191) (-1.758,-1.163) (-1.736,-1.135) (-1.714,-1.108) (-1.692,-1.081) (-1.669,-1.054) (-1.646,-1.027) (-1.623,-1.001) (-1.600,-0.975) (-1.577,-0.948) (-1.553,-0.923) (-1.529,-0.897) (-1.505,-0.871) (-1.480,-0.846) (-1.456,-0.821) (-1.431,-0.796) (-1.406,-0.772) (-1.380,-0.747) (-1.355,-0.723) (-1.329,-0.699) (-1.303,-0.676) (-1.277,-0.652) (-1.250,-0.629) (-1.224,-0.607) (-1.197,-0.584) (-1.170,-0.562) (-1.142,-0.540) (-1.115,-0.518) (-1.087,-0.496) (-1.059,-0.475) (-1.031,-0.454) (-1.002,-0.434) (-0.974,-0.413) (-0.945,-0.393) (-0.913,-0.380) (-0.878,-0.380) (-0.843,-0.380) (-0.808,-0.380) (-0.773,-0.380) (-0.738,-0.380) (-0.703,-0.380) (-0.667,-0.380) (-0.632,-0.380) (-0.597,-0.380) (-0.562,-0.380) (-0.527,-0.380) (-0.492,-0.380) (-0.457,-0.380) (-0.422,-0.380) (-0.386,-0.380) (-0.351,-0.380) (-0.316,-0.380) (-0.281,-0.380) (-0.246,-0.380) (-0.211,-0.380) (-0.176,-0.380) (-0.141,-0.380) (-0.105,-0.380) (-0.070,-0.380) (-0.035,-0.380) (-0.000,-0.380) (0.035,-0.380) (0.070,-0.380) (0.105,-0.380) (0.141,-0.380) (0.176,-0.380) (0.211,-0.380) (0.246,-0.380) (0.281,-0.380) (0.316,-0.380) (0.351,-0.380) (0.386,-0.380) (0.422,-0.380) (0.457,-0.380) (0.492,-0.380) (0.527,-0.380) (0.562,-0.380) (0.597,-0.380) (0.632,-0.380) (0.667,-0.380) (0.703,-0.380) (0.738,-0.380) (0.773,-0.380) (0.808,-0.380) (0.843,-0.380) (0.878,-0.380) (0.913,-0.380) (0.945,-0.393) (0.974,-0.413) (1.002,-0.434) (1.031,-0.454) (1.059,-0.475) (1.087,-0.496) (1.115,-0.518) (1.142,-0.540) (1.170,-0.562) (1.197,-0.584) (1.224,-0.607) (1.250,-0.629) (1.277,-0.652) (1.303,-0.676) (1.329,-0.699) (1.355,-0.723) (1.380,-0.747) (1.406,-0.772) (1.431,-0.796) (1.456,-0.821) (1.480,-0.846) (1.505,-0.871) (1.529,-0.897) (1.553,-0.923) (1.577,-0.948) (1.600,-0.975) (1.623,-1.001) (1.646,-1.027) (1.669,-1.054) (1.692,-1.081) (1.714,-1.108) (1.736,-1.135) (1.758,-1.163) (1.780,-1.191) (1.801,-1.219) (1.822,-1.247) (1.843,-1.275) (1.864,-1.303) (1.884,-1.332) (1.904,-1.361) (1.924,-1.390) (1.944,-1.419) (1.963,-1.448) (1.983,-1.477) (2.002,-1.507) (2.020,-1.537) (2.039,-1.567) (2.057,-1.597) (2.075,-1.627) (2.093,-1.657) (2.110,-1.688) (2.127,-1.718) (2.144,-1.749) (2.161,-1.780) (2.177,-1.811) (2.193,-1.842) (2.209,-1.873) (2.225,-1.905) (2.240,-1.936) (2.256,-1.968) (2.270,-2.000)};
\addplot[lineaD, flechaen=0.399] coordinates {(-1.423,-2.000) (-1.423,-1.965) (-1.422,-1.930) (-1.420,-1.895) (-1.418,-1.860) (-1.415,-1.825) (-1.412,-1.790) (-1.408,-1.755) (-1.404,-1.720) (-1.399,-1.685) (-1.393,-1.651) (-1.387,-1.616) (-1.380,-1.582) (-1.373,-1.547) (-1.366,-1.513) (-1.357,-1.479) (-1.349,-1.445) (-1.339,-1.411) (-1.330,-1.378) (-1.319,-1.344) (-1.309,-1.311) (-1.297,-1.277) (-1.285,-1.244) (-1.273,-1.212) (-1.260,-1.179) (-1.247,-1.146) (-1.233,-1.114) (-1.219,-1.082) (-1.204,-1.050) (-1.189,-1.019) (-1.173,-0.987) (-1.157,-0.956) (-1.140,-0.925) (-1.123,-0.895) (-1.105,-0.864) (-1.087,-0.834) (-1.069,-0.805) (-1.050,-0.775) (-1.030,-0.746) (-1.010,-0.717) (-0.990,-0.689) (-0.969,-0.660) (-0.948,-0.632) (-0.926,-0.605) (-0.904,-0.578) (-0.881,-0.551) (-0.858,-0.524) (-0.825,-0.520) (-0.789,-0.520) (-0.754,-0.520) (-0.719,-0.520) (-0.684,-0.520) (-0.649,-0.520) (-0.614,-0.520) (-0.579,-0.520) (-0.544,-0.520) (-0.509,-0.520) (-0.474,-0.520) (-0.439,-0.520) (-0.404,-0.520) (-0.368,-0.520) (-0.333,-0.520) (-0.298,-0.520) (-0.263,-0.520) (-0.228,-0.520) (-0.193,-0.520) (-0.158,-0.520) (-0.123,-0.520) (-0.088,-0.520) (-0.053,-0.520) (-0.018,-0.520) (0.018,-0.520) (0.053,-0.520) (0.088,-0.520) (0.123,-0.520) (0.158,-0.520) (0.193,-0.520) (0.228,-0.520) (0.263,-0.520) (0.298,-0.520) (0.333,-0.520) (0.368,-0.520) (0.404,-0.520) (0.439,-0.520) (0.474,-0.520) (0.509,-0.520) (0.544,-0.520) (0.579,-0.520) (0.614,-0.520) (0.649,-0.520) (0.684,-0.520) (0.719,-0.520) (0.754,-0.520) (0.789,-0.520) (0.825,-0.520) (0.858,-0.524) (0.881,-0.551) (0.904,-0.578) (0.926,-0.605) (0.948,-0.632) (0.969,-0.660) (0.990,-0.689) (1.010,-0.717) (1.030,-0.746) (1.050,-0.775) (1.069,-0.805) (1.087,-0.834) (1.105,-0.864) (1.123,-0.895) (1.140,-0.925) (1.157,-0.956) (1.173,-0.987) (1.189,-1.019) (1.204,-1.050) (1.219,-1.082) (1.233,-1.114) (1.247,-1.146) (1.260,-1.179) (1.273,-1.212) (1.285,-1.244) (1.297,-1.277) (1.309,-1.311) (1.319,-1.344) (1.330,-1.378) (1.339,-1.411) (1.349,-1.445) (1.357,-1.479) (1.366,-1.513) (1.373,-1.547) (1.380,-1.582) (1.387,-1.616) (1.393,-1.651) (1.399,-1.685) (1.404,-1.720) (1.408,-1.755) (1.412,-1.790) (1.415,-1.825) (1.418,-1.860) (1.420,-1.895) (1.422,-1.930) (1.423,-1.965) (1.423,-2.000)};
\addplot[lineaD, flechaen=0.614] coordinates {(-0.754,-2.000) (-0.772,-1.970) (-0.790,-1.940) (-0.806,-1.909) (-0.821,-1.877) (-0.836,-1.845) (-0.849,-1.813) (-0.862,-1.780) (-0.873,-1.747) (-0.884,-1.714) (-0.894,-1.680) (-0.902,-1.646) (-0.910,-1.612) (-0.917,-1.578) (-0.923,-1.543) (-0.928,-1.508) (-0.932,-1.474) (-0.935,-1.439) (-0.938,-1.404) (-0.939,-1.369) (-0.940,-1.334) (-0.939,-1.299) (-0.938,-1.264) (-0.936,-1.229) (-0.933,-1.194) (-0.930,-1.159) (-0.925,-1.124) (-0.920,-1.090) (-0.913,-1.055) (-0.906,-1.021) (-0.899,-0.987) (-0.890,-0.953) (-0.881,-0.919) (-0.870,-0.885) (-0.859,-0.852) (-0.848,-0.819) (-0.835,-0.786) (-0.822,-0.754) (-0.808,-0.722) (-0.793,-0.690) (-0.778,-0.658) (-0.753,-0.640) (-0.718,-0.640) (-0.683,-0.640) (-0.648,-0.640) (-0.613,-0.640) (-0.578,-0.640) (-0.543,-0.640) (-0.508,-0.640) (-0.473,-0.640) (-0.438,-0.640) (-0.403,-0.640) (-0.368,-0.640) (-0.333,-0.640) (-0.298,-0.640) (-0.263,-0.640) (-0.228,-0.640) (-0.193,-0.640) (-0.158,-0.640) (-0.123,-0.640) (-0.088,-0.640) (-0.053,-0.640) (-0.018,-0.640) (0.018,-0.640) (0.053,-0.640) (0.088,-0.640) (0.123,-0.640) (0.158,-0.640) (0.193,-0.640) (0.228,-0.640) (0.263,-0.640) (0.298,-0.640) (0.333,-0.640) (0.368,-0.640) (0.403,-0.640) (0.438,-0.640) (0.473,-0.640) (0.508,-0.640) (0.543,-0.640) (0.578,-0.640) (0.613,-0.640) (0.648,-0.640) (0.683,-0.640) (0.718,-0.640) (0.753,-0.640) (0.778,-0.658) (0.793,-0.690) (0.808,-0.722) (0.822,-0.754) (0.835,-0.786) (0.848,-0.819) (0.859,-0.852) (0.870,-0.885) (0.881,-0.919) (0.890,-0.953) (0.899,-0.987) (0.906,-1.021) (0.913,-1.055) (0.920,-1.090) (0.925,-1.124) (0.930,-1.159) (0.933,-1.194) (0.936,-1.229) (0.938,-1.264) (0.939,-1.299) (0.940,-1.334) (0.939,-1.369) (0.938,-1.404) (0.935,-1.439) (0.932,-1.474) (0.928,-1.508) (0.923,-1.543) (0.917,-1.578) (0.910,-1.612) (0.902,-1.646) (0.894,-1.680) (0.884,-1.714) (0.873,-1.747) (0.862,-1.780) (0.849,-1.813) (0.836,-1.845) (0.821,-1.877) (0.806,-1.909) (0.790,-1.940) (0.772,-1.970) (0.754,-2.000)};
\addplot[lineaD, flechaen=0.500] coordinates {(-0.657,-0.750) (-0.621,-0.750) (-0.586,-0.750) (-0.551,-0.750) (-0.516,-0.750) (-0.481,-0.750) (-0.446,-0.750) (-0.410,-0.750) (-0.375,-0.750) (-0.340,-0.750) (-0.305,-0.750) (-0.270,-0.750) (-0.234,-0.750) (-0.199,-0.750) (-0.164,-0.750) (-0.129,-0.750) (-0.094,-0.750) (-0.059,-0.750) (-0.023,-0.750) (0.012,-0.750) (0.047,-0.750) (0.082,-0.750) (0.117,-0.750) (0.152,-0.750) (0.188,-0.750) (0.223,-0.750) (0.258,-0.750) (0.293,-0.750) (0.328,-0.750) (0.363,-0.750) (0.399,-0.750) (0.434,-0.750) (0.469,-0.750) (0.504,-0.750) (0.539,-0.750) (0.574,-0.750) (0.610,-0.750) (0.645,-0.750) (0.665,-0.770) (0.671,-0.804) (0.676,-0.839) (0.680,-0.874) (0.683,-0.909) (0.684,-0.944) (0.684,-0.979) (0.683,-1.015) (0.681,-1.050) (0.677,-1.085) (0.673,-1.119) (0.667,-1.154) (0.660,-1.189) (0.651,-1.223) (0.641,-1.256) (0.630,-1.290) (0.617,-1.323) (0.604,-1.355) (0.589,-1.387) (0.572,-1.418) (0.554,-1.448) (0.535,-1.478) (0.515,-1.506) (0.493,-1.534) (0.470,-1.561) (0.446,-1.586) (0.420,-1.610) (0.394,-1.633) (0.366,-1.655) (0.337,-1.675) (0.307,-1.693) (0.276,-1.710) (0.245,-1.725) (0.212,-1.739) (0.179,-1.750) (0.145,-1.760) (0.111,-1.767) (0.076,-1.773) (0.041,-1.776) (0.006,-1.778) (-0.029,-1.777) (-0.064,-1.774) (-0.099,-1.769) (-0.134,-1.763) (-0.168,-1.754) (-0.201,-1.743) (-0.234,-1.730) (-0.266,-1.716) (-0.297,-1.699) (-0.327,-1.681) (-0.356,-1.662) (-0.385,-1.641) (-0.412,-1.618) (-0.437,-1.594) (-0.462,-1.569) (-0.486,-1.543) (-0.508,-1.516) (-0.529,-1.487) (-0.548,-1.458) (-0.566,-1.428) (-0.583,-1.397) (-0.599,-1.366) (-0.613,-1.334) (-0.626,-1.301) (-0.638,-1.268) (-0.648,-1.234) (-0.657,-1.200) (-0.664,-1.166) (-0.671,-1.131) (-0.676,-1.096) (-0.680,-1.061) (-0.683,-1.026) (-0.684,-0.991) (-0.684,-0.956) (-0.683,-0.921) (-0.681,-0.886) (-0.678,-0.851) (-0.673,-0.816) (-0.668,-0.781) (-0.657,-0.750)};
\addplot[lineaD, flechaen=0.382] coordinates {(-0.523,-0.850) (-0.488,-0.850) (-0.453,-0.850) (-0.417,-0.850) (-0.382,-0.850) (-0.346,-0.850) (-0.311,-0.850) (-0.276,-0.850) (-0.240,-0.850) (-0.205,-0.850) (-0.170,-0.850) (-0.134,-0.850) (-0.099,-0.850) (-0.063,-0.850) (-0.028,-0.850) (0.007,-0.850) (0.043,-0.850) (0.078,-0.850) (0.114,-0.850) (0.149,-0.850) (0.184,-0.850) (0.220,-0.850) (0.255,-0.850) (0.291,-0.850) (0.326,-0.850) (0.361,-0.850) (0.397,-0.850) (0.432,-0.850) (0.468,-0.850) (0.503,-0.850) (0.525,-0.864) (0.519,-0.898) (0.512,-0.933) (0.502,-0.967) (0.492,-1.001) (0.479,-1.034) (0.465,-1.066) (0.449,-1.098) (0.431,-1.129) (0.411,-1.158) (0.390,-1.186) (0.367,-1.213) (0.343,-1.239) (0.317,-1.263) (0.289,-1.285) (0.260,-1.305) (0.230,-1.324) (0.198,-1.340) (0.166,-1.354) (0.132,-1.365) (0.098,-1.374) (0.063,-1.380) (0.028,-1.383) (-0.007,-1.384) (-0.043,-1.382) (-0.078,-1.377) (-0.113,-1.370) (-0.147,-1.360) (-0.180,-1.348) (-0.212,-1.333) (-0.243,-1.316) (-0.273,-1.297) (-0.301,-1.276) (-0.328,-1.253) (-0.353,-1.228) (-0.377,-1.202) (-0.399,-1.175) (-0.420,-1.146) (-0.439,-1.116) (-0.456,-1.085) (-0.471,-1.053) (-0.485,-1.020) (-0.496,-0.987) (-0.507,-0.953) (-0.515,-0.918) (-0.522,-0.884) (-0.523,-0.850)};
\addplot[lineaD, flechaen=0.500] coordinates {(-0.338,-0.941) (-0.303,-0.940) (-0.267,-0.940) (-0.232,-0.940) (-0.196,-0.940) (-0.161,-0.940) (-0.125,-0.940) (-0.090,-0.940) (-0.054,-0.940) (-0.019,-0.940) (0.017,-0.940) (0.052,-0.940) (0.088,-0.940) (0.123,-0.940) (0.159,-0.940) (0.194,-0.940) (0.230,-0.940) (0.266,-0.940) (0.301,-0.940) (0.337,-0.940) (0.321,-0.967) (0.299,-0.994) (0.274,-1.020) (0.247,-1.043) (0.218,-1.064) (0.188,-1.082) (0.156,-1.098) (0.123,-1.111) (0.089,-1.121) (0.054,-1.128) (0.019,-1.131) (-0.017,-1.131) (-0.052,-1.128) (-0.087,-1.122) (-0.121,-1.112) (-0.154,-1.099) (-0.186,-1.083) (-0.217,-1.065) (-0.245,-1.044) (-0.272,-1.021) (-0.297,-0.996) (-0.320,-0.969) (-0.338,-0.941)};
\addplot[lineaD, flechaen=0.087, flechaen=0.913] coordinates {(-2.600,0.000) (2.600,0.000)};

  \draw[texto, line width=0.6pt] (axis cs:0,0) circle[radius=1];
  \draw[vecP] (axis cs:-0.45,-0.55) -- (axis cs:0.45,-0.55);
  \node[etiqm, font=\footnotesize] at (axis cs:0,-0.78) {$\vP$};
\end{axis}
\end{tikzpicture}
\caption{Líneas de $\vD$}
\end{subfigure}
\caption{Esfera uniformemente polarizada (\cref{prob:esfera-polarizada}); trazos: equipotenciales. (a) Fuera,
el campo es exactamente el de un dipolo $\vp=\tfrac43\pi a^3\vP$ en el centro; dentro es uniforme y
opuesto a la polarización, $\vE=-\vP/3\eo$ (campo despolarizante): las líneas nacen en la carga ligada
positiva $\sigmaP=P\cos\theta$ y mueren en la negativa. (b) Sin carga libre, $\vD$ es solenoidal y sus
líneas son cerradas; dentro, $\vD=\eo\vE+\vP=\tfrac23\vP$ tiene el sentido de $\vP$.}
\label{fig:esfera-polarizada}
\end{figure}
